\documentclass[11pt]{article}

% ============================================================================
% Apparatus axis (Paper A) — modular statements-of-record root.
%
% main.tex carries the preamble + front matter + the ordered \input of the
% sections/ and appendices/ files (the labeled environments live there; the
% inventory extractor reads those files). Preamble adopted from
% ../paper-template (publication style: lmodern/microtype, titling, fancyhdr
% footer, orcidlink author, caption/booktabs); BSD macros are in
% includes/paper_macros.tex.
%
% Label schema: <kind>:apparatus:<short-name>, kind-prefix matching env:
%   definition->def  theorem->thm  remark->rmk  warning->warn
%   obligation->obl  nonclaim->nonclaim
% ============================================================================

\usepackage[T1]{fontenc}
\usepackage[utf8]{inputenc}
\usepackage[english]{babel}
\usepackage[margin=1in]{geometry}
\usepackage{lmodern}
\usepackage{microtype}
\usepackage{mathtools}
\usepackage{amsmath,amssymb,amsthm}
\usepackage{array}
\usepackage{booktabs}
\usepackage{tabularx}
\usepackage{longtable}
\usepackage{float}
\usepackage{graphicx}
\usepackage{tikz}
\usetikzlibrary{arrows.meta,calc,fit,positioning}
\usepackage{pdflscape}
\usepackage{caption}
\usepackage{enumitem}
\usepackage{fancyhdr}
\usepackage[numbers]{natbib}
\usepackage{doi}
\usepackage{xurl}
\usepackage{hyperref}
\usepackage{cleveref}
\usepackage{orcidlink}
\usepackage{titling}

\usepackage{aliascnt}
\newtheorem{theorem}{Theorem}[section]
% Shared-counter environments get alias counters so that cleveref prints the
% environment name (Definition, Obligation, ...) rather than "Theorem".
\newcommand{\sharedtheorem}[2]{%
  \newaliascnt{#1}{theorem}%
  \newtheorem{#1}[#1]{#2}%
  \aliascntresetthe{#1}}
\sharedtheorem{lemma}{Lemma}
\sharedtheorem{proposition}{Proposition}
\sharedtheorem{corollary}{Corollary}
\theoremstyle{definition}
\sharedtheorem{definition}{Definition}
\sharedtheorem{example}{Example}
\sharedtheorem{remark}{Remark}
\sharedtheorem{warning}{Warning}
\sharedtheorem{obligation}{Obligation}
\sharedtheorem{nonclaim}{Nonclaim}

\urlstyle{same}
\captionsetup{font=small, labelfont=bf, labelsep=period}
\captionsetup[table]{skip=6pt, position=top}
\captionsetup[figure]{skip=6pt}
\crefname{theorem}{theorem}{theorems}
\Crefname{theorem}{Theorem}{Theorems}
\crefname{lemma}{lemma}{lemmas}
\Crefname{lemma}{Lemma}{Lemmas}
\crefname{proposition}{proposition}{propositions}
\Crefname{proposition}{Proposition}{Propositions}
\crefname{corollary}{corollary}{corollaries}
\Crefname{corollary}{Corollary}{Corollaries}
\crefname{definition}{definition}{definitions}
\Crefname{definition}{Definition}{Definitions}
\crefname{remark}{remark}{remarks}
\Crefname{remark}{Remark}{Remarks}
\crefname{warning}{warning}{warnings}
\Crefname{warning}{Warning}{Warnings}
\crefname{obligation}{obligation}{obligations}
\Crefname{obligation}{Obligation}{Obligations}
\crefname{nonclaim}{nonclaim}{nonclaims}
\Crefname{nonclaim}{Nonclaim}{Nonclaims}
\crefname{example}{example}{examples}
\Crefname{example}{Example}{Examples}
\crefname{appendix}{appendix}{appendices}
\Crefname{appendix}{Appendix}{Appendices}
\setlist[enumerate]{leftmargin=*, itemsep=2pt, topsep=4pt}
\setlength{\droptitle}{-3.2em}
\pretitle{\begin{center}\LARGE\bfseries}
\posttitle{\par\end{center}\vspace{0.5em}}
\preauthor{\begin{center}\normalsize}
\postauthor{\par\end{center}\vspace{-0.4em}}
\predate{\begin{center}\small}
\postdate{\par\end{center}\vspace{0.6em}}
\clubpenalty=10000
\widowpenalty=10000
\displaywidowpenalty=10000
\interfootnotelinepenalty=10000
\allowdisplaybreaks[2]
\fancypagestyle{firstpage}{\fancyhf{}
  \fancyfoot[C]{\scriptsize Automorph Inc., Wilmington, DE, USA \quad \texttt{ioannis@automorph.io} \quad \textcopyright\ 2026 \quad CC-BY 4.0 \quad \doi{10.5281/zenodo.20713981}}
  \renewcommand{\headrulewidth}{0pt}
  \renewcommand{\footrulewidth}{0pt}
}

% BSD macros + the math-heavy-skeleton layout settings (\emergencystretch, \sloppy).
% Paper macros for the apparatus axis (Phase 1). \input by main.tex.
% Reserved-symbol policy binds Lean naming (see paper/notation_and_terminology.md).
% This file is the canonical macro source; paper/notes/macro-audit.md mirrors it.

% --- Shared / framework macros (identical block in both papers) ---
\newcommand{\Q}{\mathbb{Q}}
\newcommand{\R}{\mathbb{R}}
\newcommand{\Z}{\mathbb{Z}}
\newcommand{\Qp}{\mathbb{Q}_p}
\newcommand{\Zp}{\mathbb{Z}_p}
\newcommand{\Sha}{\operatorname{Sha}}
\newcommand{\Reg}{\operatorname{Reg}}
\newcommand{\Tam}{\operatorname{Tam}}
\newcommand{\tors}{\operatorname{tors}}
\newcommand{\tr}{\operatorname{tr}}
\newcommand{\Pf}{\operatorname{Pf}}
\newcommand{\Fitt}{\operatorname{Fitt}}
\newcommand{\cores}{\operatorname{cores}}
\newcommand{\rank}{\operatorname{rank}}
\newcommand{\Hf}{H^1_f}
\newcommand{\kapr}{\kappa_r}

% --- Apparatus-specific macros ---
\newcommand{\XiC}{\Xi_C(D\mid L)}        % adequacy / Schur residual
\newcommand{\KDD}{K_{DD}}
\newcommand{\KDL}{K_{DL}}
\newcommand{\KLL}{K_{LL}}
\newcommand{\Kdagger}{K_{LL}^{\dagger}}  % pseudo-inverse block
\newcommand{\vsrc}{\mathrm{vsrc}}        % exterior-algebra calibrator
\newcommand{\Jone}{J_1}
\newcommand{\Jtwo}{J_2}
\newcommand{\Jtwelve}{J_{12}}
\newcommand{\MCphi}{M_C^{\varphi}}       % typed column (Lean M_C_phi)

% Math-heavy SoR skeleton: allow interword stretch to avoid overfull boxes
\setlength{\emergencystretch}{3.5em}
\sloppy


\title{Decomposing the Bloch--Kato Fundamental Line: Adequacy and No-Go
Results for the BSD Invariants}
\author{Ioannis Tsiokos\,\orcidlink{0009-0009-7659-5964}}
\date{v1: 31 August 2026\quad\textperiodcentered\quad v2: 3 October 2026}
\hypersetup{
  pdftitle={Decomposing the Bloch--Kato Fundamental Line: Adequacy and No-Go Results for the BSD Invariants},
  pdfauthor={Ioannis Tsiokos},
  hidelinks
}

\begin{document}
\maketitle
\thispagestyle{firstpage}

\begin{abstract}
The Birch--Swinnerton-Dyer formula for an elliptic curve $E/\Q$ is a
product of pieces of different kinds: a period, a regulator, Tamagawa
numbers, torsion and the order of the Tate--Shafarevich group, which in
the Bloch--Kato formulation are joined with local $p$-adic data and
orientation conventions in one determinant line.  Each piece is the image
of richer data.  We separate these data into five columns
(height/regulator, analytic/period, finite source, local $p$-adic, and
determinant assembly) and ask, column by column, which data a given
number remembers.  The test is a finite Schur complement, which vanishes
exactly when the desired readout is a linear function of the admitted
data.  The full data of each of the first four columns are adequate for
their readouts, and the five columns together are adequate for an
orientation-normalized Bloch--Kato scalar, which forgets the orientation.
Along the way we fix three normalization points: the invariant finite
factor is the order of a paired group, not a numerical Pfaffian of its
Cassels--Tate pairing; on the finite local condition the relevant map is
the logarithm, not the dual exponential; and the assembled value must not
be reused as a factor.  We then identify readouts that lose information.
The most concrete is arithmetic: the curves $571a1$ and $2045b1$ have the
same rank, Tamagawa numbers, torsion and $\dim\Sha[2]$, but $2$-primary
Tate--Shafarevich groups of orders $4$ and $16$.  Other results, proved on
explicit models, concern the scalar finite factor, Gram matrices in rank
at least two, and finite lists of curves and of primes.  We also identify
the normalization $\kapr=\#\Sha/\#E(\Q)_{\tors}^{\,2}$ under which a
cascade form of the leading term agrees with the standard one.  The
comparisons from the columns to the Bloch--Kato line are stated as
external inputs.  Nothing here proves BSD.  Many statements are checked
in Lean~4, and each is labelled by the exact extent of that check.
\end{abstract}

\medskip
\noindent\textbf{Keywords.} Birch--Swinnerton-Dyer conjecture; Bloch--Kato fundamental line; elliptic curves; Tate--Shafarevich group; Cassels--Tate pairing; Schur complement; regulator; formal verification (Lean); Six Birds Theory.

\medskip
\noindent
\textbf{Conventions.} $E/\Q$ is an elliptic curve, $M_E=h^1(E)(1)$,
$r=\rank E(\Q)$ and $r_{\mathrm{an}}=\operatorname{ord}_{s=1}L(E,s)$.
Curve labels are Cremona labels.

% ===========================================================================
\section{Introduction}\label{sec:introduction}
% statements-of-record rows resolved here: none

The Birch--Swinnerton-Dyer formula~\citep{BirchSwinnertonDyer1965,Tate1966bsd} is usually first met as a scalar
identity.  For an elliptic curve $E/\Q$ of rank $r$, the standard leading
coefficient form compares
\[
  \frac{L^{(r)}(E,1)}{r!}
  \quad\text{with}\quad
  \Omega_E\,\Reg(E/\Q)\,
  \frac{\#\Sha(E/\Q)\,\Tam(E)}{\#E(\Q)_{\tors}^{\,2}} .
\]
This product is already visibly composite: an analytic leading term and
period, a height regulator, a finite Tamagawa--torsion--$\Sha$ factor, local
$p$-adic data, and the determinant-line conventions that make all of these
terms live in a common orientation.  The question addressed in this paper is
not whether the scalar BSD identity is true.  It is the more local diagnostic
question: when one separates the Bloch--Kato fundamental line~\citep{BlochKato1990,FontainePerrinRiou1994} into these
typed pieces, which pieces carry enough native data to recover their intended
readout, and where is information lost by passing to a public scalar shadow?

The word ``typed'' is used here in its elementary sense.  We keep the five
classes of data in five labelled slots rather than immediately multiplying
them into one number.  The height/regulator slot contains the Mordell--Weil
lattice and the Neron--Tate pairing; the analytic/period slot contains the
leading Taylor coefficient and the real period; the finite slot contains
Tamagawa factors, torsion, the $p$-primary $\Sha$ data, and Cassels--Tate
pairings; the $p$-adic slot contains the ordinary or signed local data; and
the determinant slot contains the orientation, sign, and unit conventions
needed for a single determinant line.  This decomposition is bookkeeping
only in the innocuous sense that a direct-sum decomposition is: it is
precisely the bookkeeping needed to ask whether each column is adequate.

The adequacy test is a Schur-complement calculation~\citep{Zhang2005schur,HornJohnson2013}.  Suppose that $L$ is
the native family admitted in a column and $D$ is the readout one wants to
recover from that column.  Form the block Gram quantities
\[
  \KLL,\qquad \KDL,\qquad \KDD,
\]
where $\KLL$ records the covariance of the admitted source directions,
$\KDL$ records the pairing between the readout and those source directions,
and $\KDD$ records the readout's own pre-source energy.  The adequacy
residual is
\[
  \XiC=\KDD-\KDL\,\Kdagger\,K_{LD},
\]
with $\Kdagger$ the Moore--Penrose inverse of $\KLL$ in the finite
calculation at hand.  The column is called \emph{adequate} when
$\XiC=0$: after projecting the desired readout onto the admitted source
directions, no residual direction remains.  If $\XiC>0$ in a scalar trace
calculation, then the admitted row is not enough; it has forgotten a
direction that is present in the typed target.

All adequacy computations in this paper are finite and explicit.  The
linear algebra is deliberately elementary: finite index sets, row vectors,
dot products, identity matrices, matrix multiplication, and traces.  The
collapses are not appeals to a black-box functional-analytic statement.
They are the finite identity-matrix action
\[
  I v=v
\]
and its Schur-complement consequence.  For example, if the full native
height matrix is admitted, the readout row $d_H$ for
$d(\log\det)_H$ projects through the identity source, so
$d_Hd_H^{\mathsf T}-d_H I d_H^{\mathsf T}=0$.  If only a scalar determinant
is admitted, the orthogonal complement has positive trace.  The same
pattern recurs in the analytic/period, finite-source, ordinary $p$-adic,
signed $p$-adic, and rank-two height calculations.

This finite Schur-complement language is the specialization, for the BSD
fundamental line, of the adequacy calculus developed in the surrounding
Six Birds framework \citep{TsiokosNeedles2026,TsiokosFoundationsII2026}.
For the reader of the present paper, no framework machinery is needed
beyond the following translation.  An ``adequacy residual'' is simply the
Schur complement above.  A ``typed five-column decomposition'' is simply
the decision to keep the height, period, finite, $p$-adic, and determinant
data in separate labelled columns.  An ``audit-result no-go'' is a recorded
boundary statement: it identifies where the theorem-grade literature, in
the audited form used here, does not supply the missing identity.  The
interaction vocabulary and audit discipline are inherited from the finite
calculus of \citet{TsiokosFoundationsIII2026}; the mathematical statements
below are written directly in the language of the BSD factors.

A low-conductor example makes the separation concrete.  For the curve
$11a1$~\citep{Cremona1997,LMFDB}, the scalar rank-zero BSD comparison has
\[
  \frac{L(E,1)}{\Omega_E^+}=0.2
  \qquad\text{and}\qquad
  \frac{\prod_v c_v\,\#\Sha(E/\Q)}{\#E(\Q)_{\tors}^{\,2}}
  =\frac{5}{25}=0.2.
\]
As a scalar check this is a single equality.  In the five-column atlas it is
not a single source.  The analytic/period column contributes the quotient
$L(E,1)/\Omega_E^+$; the height/regulator column is rank-zero and therefore
vacuous; the finite column remembers that the number $5/25$ arose from a
Tamagawa factor and torsion rather than from Cassels--Tate pairing data;
the local $p$-adic columns ask which ordinary or signed source data are
present at each prime; and the determinant column records the orientation
needed to assemble the product.  The scalar equality is therefore compatible
with, but coarser than, the typed comparison.  The point of the apparatus is
to measure this coarsening column by column.

This is also the reason for working with readouts rather than only with
equalities of final products.  A scalar product can be correct for several
different reasons.  In the finite factor, a valuation can sit in the
Tamagawa part or in the $\Sha$ part and still contribute the same exponent
to the public scalar quotient.  In the signed supersingular local theory,
the two signed directions can have the same unsigned aggregate while
differing in their $+$ and $-$ components.  In rank at least two, a Gram
matrix can give the regulator once a Mordell--Weil basis is already known,
but the matrix itself does not name independent motivic-cycle sources.
These are not paradoxes.  They are exactly what one should expect when a
multi-column object is projected to a single row.  The Schur residual is the
linear-algebraic device that records which directions survive the projection
and which directions have been lost.

The paper therefore uses ``native family'' in a modest sense: it means the
data one is willing to admit as the source for a column before applying the
readout.  For the analytic/period column, the native family is the pair
$(A_E,\Omega_E^+)$, where
$A_E=L^{(r_{\mathrm{an}})}(E,1)/r_{\mathrm{an}}!$ is the analytic leading
coefficient and $\Omega_E^+$ is the real period, rather than the quotient
alone.  For the finite column, it is the tower containing Tamagawa, torsion,
$\Sha$, and Cassels--Tate
pairing data rather than the scalar finite factor alone.  For the signed
$p$-adic column, it is the signed pair of local directions rather than their
unsigned sum.  The adequacy theorem for a column says that this admitted
native family is large enough for the intended readout.  A no-go theorem
says that a smaller public row is not large enough, and it identifies an
explicit pair of inputs or an explicit residual direction witnessing the
loss.

Nothing in this setup requires the reader to adopt the surrounding
terminology as a black box.  The relevant matrices are small: identity
matrices of sizes $2$, $3$, $4$, and $5$, rank-one projectors, and diagonal
traces.  The Moore--Penrose notation is used only where the finite matrix
has an evident inverse or pseudoinverse in the displayed calculation.  The
rank-one residuals are ordinary rational computations such as
\[
  \tr(I_4-P_s)=4-1=3
\]
for the scalar finite row $s=(1,-2,1,0)$, or
\[
  \tr(I_3-P_{\mathrm{uns}})=3-2=1
\]
for the unsigned signed-local projection.  Thus the reader can check the
adequacy and inadequacy claims at the same level at which the statements are
made: finite-dimensional linear algebra over explicit coordinates.

The first part of the paper constructs the five columns.  The
Bloch--Kato component map sends the integrated BSD carrier to the
determinant-line target by the direct sum of the height/regulator,
analytic/period, finite, local $p$-adic, and determinant-assembly
components.  After choosing a comparison trivialization, the corresponding
defect decomposes additively into the five component defects.  Consequently,
if every component defect vanishes in the declared comparison, then the
total Bloch--Kato defect vanishes; contrapositively, a nonzero total defect
must be witnessed by at least one nonzero or legally incomparable component.

The order of the paper follows this diagnostic structure.  The
decomposition section fixes the five columns and the additive defect
equation.  The next four sections treat the height/regulator,
analytic/period, finite-source, and $p$-adic columns separately, always
asking the same question: does the declared native family collapse the
Schur residual for the declared readout?  The determinant-assembly section
then records the one remaining orientation coordinate needed to assemble
the local tensor into a global determinant-line object.  The comparison
section contracts the typed atlas back to the usual Bloch--Kato scalar
readout and proves that the typed source is adequate for that contraction.

The middle sections prove the adequacy collapses.  The height/regulator
column admits the full Neron--Tate pairing and collapses by the identity
action.  The analytic/period column admits both logarithmic coordinates
$(\log A_E,\log\Omega_E^+)$ and collapses for the readout
$\log A_E-\log\Omega_E^+$.  The finite-source column collapses only when
the Tamagawa, torsion, $\Sha$, and Cassels--Tate coordinates are all present;
the scalar finite factor and the dimension-only $\Sha$ row leave positive
residuals.  The ordinary $p$-adic row collapses with ordinary-control data
included, while omitting that coordinate leaves a one-dimensional residual.
The signed supersingular row collapses only after separating the $+$ and
$-$ directions; unsigned aggregation forgets their difference.  The
determinant-assembly calculation then shows a structural residual exactly
in the determinant-orientation coordinate when the four constructed
components are assembled without the final determinant row.

The comparison section contracts the typed five-column tensor to the usual
Bloch--Kato scalar package.  At the log-tangent level the comparison is the
row $(1,1,1,1,0)$: it reads the height, analytic, finite, and local
$p$-adic columns and kills the determinant-orientation coordinate.  Its
Schur residual is zero against the typed source, so the typed source is
adequate for the scalar Bloch--Kato readout.  At the same time, the kernel
contains the determinant coordinate, so the scalar package is a strict
one-dimensional shadow of the typed atlas.

Two calibrations complete the apparatus.  The first is the exterior algebra
$\vsrc$, generated in rank two by symbols $\Jone,\Jtwo$ and
$\Jtwelve=\Jone\Jtwo$ with $\Jone^2=\Jtwo^2=0$ and
$\Jtwo\Jone=-\Jtwelve$.  It is a determinant-shaped source algebra: the
wedge $\Jone\wedge\Jtwo$ records independence of two virtual motivic-cycle
directions without pretending that the corresponding motivic cycles have
been constructed.  The second is the normalization constant $\kapr$.  Once
the cascade convention $c_E=\Tam(E)$ is fixed, the cascade leading-term
form is equivalent to the standard Strong-BSD leading-coefficient form
exactly when
\[
  \kapr(E)=\frac{\#\Sha(E/\Q)}{\#E(\Q)_{\tors}^{\,2}}.
\]
This is an equivalence of forms, not a derivation of BSD.

The calibration sections have a different purpose from the column
sections.  The column sections test whether a given source row is adequate
for an already specified readout.  The $\vsrc$ section asks what sort of
formal source algebra would be capable of expressing rank-two independence
before any descent to actual motivic cycles is supplied.  The
$\kapr$ section asks how the cascade normalization must be chosen so that
the cascade leading-term expression is not a new conjecture but the
standard Strong-BSD expression written in the apparatus coordinates.  These
two calibrations constrain the apparatus at opposite ends: $\vsrc$ marks
the source-side independence problem, while $\kapr$ fixes the scalar
normalization at the target side.

The final mathematical contribution is the bounded no-go suite.  The
finite scalar factor is noninjective on the typed finite tower; the
dimension of $\Sha[p]$ is a shadow of Cassels--Tate determinant data; the
Neron--Tate Gram matrix does not by itself construct rank-$r$ motivic-cycle
sources; a five-curve window cannot promote to a global theorem depending
only on those five rows; and no fixed finite prime list certifies all local
support-prime rows.  These statements are useful because they are local and
typed: each identifies a precise missing direction rather than replacing
the missing direction by an undifferentiated failure of BSD.

The scope is deliberately narrow.  The no-go suite (NG-1--NG-5,
higher-rank, global, support-prime) records where theorem-grade Mode-A
literature (the established arithmetic results audited as external
provenance) does \textbf{not} reach; it is \textbf{not} a claim that the
missing identities cannot exist.  Adequacy collapses are proved at the
finite, mathlib-free, typed level (identity-matrix action; $\Q$ rank-1
projector traces); operator/motivic full generality is not claimed.  The
$\vsrc$ exterior-product relations are represented structurally on the
$2^r$ coordinate space; the dimension/calibration content is faithful,
full axiomatization is not claimed.  The five per-column transports are
recognition sources (carriers), not derived theorems.

\section{Decomposition: the typed five columns}
\label{sec:decomposition}

The introduction separated the scalar BSD product into five kinds of data:
height/regulator, analytic/period, finite-source, local $p$-adic, and
determinant-assembly data.  We now make this separation into a typed object.
The point is not to replace the Bloch--Kato determinant line by five
unrelated objects, but to keep track of the five native sources before they
are assembled into the common determinant-line target.

Fix an elliptic curve $E/\Q$ and put $M_E=h^1(E)(1)$.  The five columns are
as follows.  The height/regulator column records the Mordell--Weil lattice
and the Neron--Tate determinant.  The analytic/period column records the
leading coefficient and period data.  The finite-source column records
Tamagawa factors, torsion, the $p$-primary $\Sha$ data, and Cassels--Tate
pairings.  For each prime $p$, the local $p$-adic column records the
ordinary or signed local source data.  The determinant-assembly column
records the orientation, sign, and unit conventions needed to compare all
four arithmetic columns in one determinant line.

The five-column map and the common Schur-residual test are summarized in
\Cref{tab:five-column-map,fig:five-column-schematic}.  The table fixes the
paper-prose names used in the later column sections; the figure records the
same bookkeeping as a source-to-target diagram.
\begin{table}[tbp]
\centering
\small
\caption{The five typed columns and their adequacy residuals.  Each row
separates the native source data from the determinant-line target before the
scalar BSD product is read out.}
\label{tab:five-column-map}
\begin{tabularx}{\textwidth}{@{}>{\raggedright\arraybackslash}p{0.18\textwidth}
>{\raggedright\arraybackslash}p{0.25\textwidth}
>{\raggedright\arraybackslash}p{0.22\textwidth}
>{\raggedright\arraybackslash}X@{}}
\toprule
Column & Carrier map & Residual & Adequate when \\
\midrule
Height/regulator &
\(\rho_{\mathrm{ht/reg}}^{\R}\) from the Mordell--Weil lattice and
Neron--Tate pairing &
\(\XiC\) for \(d(\log\det)_H(\delta H)=\tr(H^{-1}\delta H)\) &
the full height pairing is admitted, so the identity source kills the Schur
residual. \\
\addlinespace
Analytic/period &
\(\rho_{\mathrm{an/period}}^{\R}\) from \((A_E,\Omega_E^+)\) &
\(\XiC\) for \(\log|A_E|-\log\Omega_E^+\) (\(A_E\neq0\)) &
both logarithmic coordinates are present; either coordinate alone leaves a
one-dimensional residual. \\
\addlinespace
Finite/Tamagawa &
\(\rho_{\mathrm{finite}}^{\R}\) from Tamagawa, torsion, \(\Sha[p^\infty]\),
and Cassels--Tate data &
\(\XiC\) for the typed finite-source readout &
the full finite-source tower, including the Cassels--Tate pairing, is
retained. \\
\addlinespace
\(p\)-adic &
\(\rho_p\) from ordinary or signed local data &
\(\XiC\) for the ordinary/signed local readout &
the appropriate ordinary or signed local source is admitted, not merely an
unsigned public shadow. \\
\addlinespace
Determinant assembly &
\(\rho_{\det}\) for orientation, sign, and unit compatibility &
\(\XiC\) for the assembled determinant comparison &
the determinant-orientation coordinate is included with the four arithmetic
columns. \\
\bottomrule
\end{tabularx}
\end{table}

\begin{figure}[tbp]
\centering
\begin{tikzpicture}[
  font=\small,
  column/.style={draw, rounded corners=2pt, align=center, minimum width=2.6cm,
    minimum height=1.05cm, fill=black!3},
  residual/.style={draw, rounded corners=2pt, align=center, minimum width=7.5cm,
    minimum height=0.95cm, fill=white},
  target/.style={draw, rounded corners=2pt, align=center, minimum width=4.1cm,
    minimum height=1.05cm, fill=black!8},
  arrow/.style={-{Latex[length=2mm]}, thick}
]
\node[column] (ht) at (-5.6,1.7) {height\\regulator};
\node[column] (an) at (-2.8,1.7) {analytic\\period};
\node[column] (fin) at (0,1.7) {finite\\Tamagawa};
\node[column] (pad) at (2.8,1.7) {\(p\)-adic\\local};
\node[column] (det) at (5.6,1.7) {determinant\\assembly};
\node[residual] (xi) at (0,0) {Schur residual:\quad
  \(\XiC=K_{DD}-K_{DL}K_{LL}^{\dagger}K_{LD}\)};
\node[target] (bk) at (0,-1.55) {Bloch--Kato determinant-line target};

\foreach \n/\shift in {ht/-3cm,an/-1.5cm,fin/0cm,pad/1.5cm,det/3cm}
  \draw[arrow] (\n.south) .. controls +(0,-0.45) and +(0,0.45) ..
    ([xshift=\shift]xi.north);

\draw[arrow] (xi.south) -- (bk.north);
\node[align=center, font=\footnotesize] at (0,-2.6)
  {adequacy means the admitted native source leaves no Schur residual};
\end{tikzpicture}
\caption{The five native source columns are tested before they are assembled in
the common determinant-line target.}
\label{fig:five-column-schematic}
\end{figure}


The Bloch--Kato target is the determinant-line package attached to $M_E$~\citep{Deligne1987}:
the motive, its finite determinant line, the zeta element, and the chosen
trivialization.  The integrated BSD carrier is the corresponding typed
source object: it keeps the height/regulator, analytic/period,
finite-source, local $p$-adic, and determinant-assembly data in their
separate slots.  The first structural map of the paper is the
component-assembly map from this typed source to the determinant-line
target.

\begin{definition}[Bloch--Kato component-assembly map]
\label{def:apparatus:bk-component-map}
For $M_E=h^1(E)(1)$, let
\[
  C_{\mathrm{BSD}}(E)
  =
  \bigl(C_{\mathrm{ht}}(E),C_{\mathrm{an}}(E),C_{\mathrm{fin}}(E),
  \{C_p(E)\}_p,C_{\det}(E)\bigr)
\]
be the integrated BSD carrier with its typed height/regulator,
analytic/period, finite-source, local $p$-adic, and
determinant-assembly slots.  The source-side determinant slot
\(C_{\det}(E)\) is distinct from the Bloch--Kato target below.  Let
\[
  C_{\mathrm{BK}}(E)
  =
  \bigl(M_E,\Delta_f(M_E),z_{\mathrm{BK}}(M_E),\tau_{\mathrm{BK}}(M_E)\bigr)
\]
be the Bloch--Kato determinant-line target.  The component-assembly map is
the direct sum
\[
  \rho_{\mathrm{BK}}
  =
  \rho_{\mathrm{ht}}
  \oplus \rho_{\mathrm{an}}
  \oplus \rho_{\mathrm{fin}}
  \oplus \Bigl(\bigoplus_p \rho_p\Bigr)
  \oplus \rho_{\det}
  \;:\;
  C_{\mathrm{BSD}}(E)\longrightarrow C_{\mathrm{BK}}(E).
\]
Here $\rho_{\mathrm{ht}}$ sends the Mordell--Weil height and regulator
data to the global cohomological determinant side; $\rho_{\mathrm{an}}$
sends central leading-term and period data to the analytic
zeta/trivialization side; $\rho_{\mathrm{fin}}$ sends Tamagawa, torsion,
$\Sha[p^\infty]$, and Cassels--Tate data to the finite arithmetic
determinant side; $\rho_p$ sends the local $p$-adic analytic/Selmer data,
control data, and local normalization to the $p$-primary determinant side;
and $\rho_{\det}$ records the orientation, sign, and unit compatibility in
the same determinant line.
\end{definition}

The definition is a decomposition of the source side, not a theorem about
the truth of BSD.  Its role is to make the columns visible enough that the
adequacy residual of each column can be tested.  Once the columns are
visible, one may compare them through a common determinant-line
trivialization.  In additive coordinates, the comparison defect splits into
one contribution from each column.

\begin{definition}[Bridge-defect equation]
\label{def:apparatus:bridge-defect-equation}
After choosing a comparison trivialization of the determinant-line target,
write $\Delta_{\mathrm{BSD}}^{\mathrm{BK}}$ for the total Bloch--Kato defect
of the BSD bridge.  The bridge-defect equation is the additive
decomposition
\[
  \Delta_{\mathrm{BSD}}^{\mathrm{BK}}
  =
  E_{\mathrm{ht/reg}}
  +E_{\mathrm{an/period}}
  +E_{\mathrm{finite}}
  +\sum_p E_p
  +E_{\det}.
\]
The five summands are respectively the height/regulator defect, the
analytic/period defect, the finite-source defect, the local $p$-adic
defects, and the determinant-assembly defect.  Equivalently, in
multiplicative coordinates the same assertion says that the total component
ratio is the product of the corresponding five component ratios.
\end{definition}

\begin{theorem}[Component-killing]
\label{thm:apparatus:component-killing}
Assume that the five component defects are all legally comparable in the
chosen determinant-line trivialization.  If
\[
  E_{\mathrm{ht/reg}}=0,\qquad
  E_{\mathrm{an/period}}=0,\qquad
  E_{\mathrm{finite}}=0,\qquad
  E_p=0\ \text{for every }p,\qquad
  E_{\det}=0,
\]
then $\Delta_{\mathrm{BSD}}^{\mathrm{BK}}=0$.  Equivalently, if
$\Delta_{\mathrm{BSD}}^{\mathrm{BK}}\neq 0$, then at least one component
defect is nonzero or one of the asserted component comparisons is not
legally available.\footnote{Verified in Lean as
\texttt{SixBirdsBSD.Apparatus.Decomposition.componentKilling}; see App~\ref{app:formalization}.}
\end{theorem}

\begin{proof}
The comparison trivialization puts all five summands in the same additive
defect group.  Hence the bridge-defect equation of
Definition~\ref{def:apparatus:bridge-defect-equation} applies:
\[
  \Delta_{\mathrm{BSD}}^{\mathrm{BK}}
  =
  E_{\mathrm{ht/reg}}
  +E_{\mathrm{an/period}}
  +E_{\mathrm{finite}}
  +\sum_p E_p
  +E_{\det}.
\]
Under the stated hypotheses every summand on the right is zero.  The sum of
zero component defects is zero, so
$\Delta_{\mathrm{BSD}}^{\mathrm{BK}}=0$.

For the contraposition, suppose
$\Delta_{\mathrm{BSD}}^{\mathrm{BK}}\neq0$.  If all five component
comparisons were legally available and every component defect vanished, the
first part of the proof would force
$\Delta_{\mathrm{BSD}}^{\mathrm{BK}}=0$, a contradiction.  Therefore either
some declared component defect is nonzero, or the attempted componentwise
comparison is not legal in the determinant-line trivialization under
discussion.
\end{proof}

\begin{theorem}[Finite scalar factor is a public shadow]
\label{thm:apparatus:finite-scalar-public-shadow}
Let
\[
  C_{\mathrm{fin}}
  =
  C_{\mathrm{Tamagawa}}\oplus C_{\tors}\oplus
  C_{\Sha[p^\infty]}\oplus C_{\mathrm{CT}}
\]
be the typed finite-source tower.  The public scalar finite factor
\[
  \phi_{\mathrm{fin}}(C_{\mathrm{fin}})
  =
  \frac{|\Sha|\prod_\ell c_\ell}{|E(\Q)_{\tors}|^2}
\]
is not injective on this typed tower.  Equivalently, at a fixed prime $p$,
the valuation readout
\[
  q_p
  =
  v_p|\Sha|+v_p\Bigl(\prod_\ell c_\ell\Bigr)
  -2v_p|E(\Q)_{\tors}|
\]
does not determine the typed finite-source datum.  Thus a scalar
finite-factor row cannot be promoted to a finite-source carrier row without
additional source data.\footnote{Verified in Lean as
\texttt{SixBirdsBSD.Apparatus.Decomposition.finiteScalarPublicShadow}; see
App~\ref{app:formalization}.}
\end{theorem}

\begin{proof}
Fix a prime $p$ and consider only the three valuation coordinates appearing
in $q_p$: the $p$-primary part of $\Sha$, the Tamagawa product, and torsion.
The typed finite tower remembers which block carries a valuation, while the
scalar readout only remembers the integer combination
\[
  v_p|\Sha|+v_p\Bigl(\prod_\ell c_\ell\Bigr)
  -2v_p|E(\Q)_{\tors}|.
\]
Now compare two typed finite inputs.  In the first, put the valuation in the
$\Sha[p^\infty]$ block:
\[
  \bigl(v_p|\Sha|,\ v_p(\prod_\ell c_\ell),\
  v_p|E(\Q)_{\tors}|\bigr)=(1,0,0).
\]
In the second, put it in the Tamagawa block:
\[
  \bigl(v_p|\Sha|,\ v_p(\prod_\ell c_\ell),\
  v_p|E(\Q)_{\tors}|\bigr)=(0,1,0).
\]
The two typed inputs are distinct, since the first records a
$p$-primary $\Sha$ contribution and the second records a Tamagawa
contribution.  Nevertheless the scalar valuation is the same:
\[
  q_p=1+0-2\cdot0=1
  \qquad\text{and}\qquad
  q_p=0+1-2\cdot0=1.
\]
Thus the scalar map identifies two different typed finite-source rows.
Moreover the Cassels--Tate datum is absent from the displayed expression
for $q_p$ altogether.  Changing the Cassels--Tate pairing while holding the
three scalar valuation coordinates fixed leaves $q_p$ unchanged.  Hence the
scalar finite factor is a noninjective public shadow of the typed finite
tower.

This is an audit no-go in the sense used throughout the paper: it records
that the scalar finite row alone does not contain enough information to
recover the typed finite-source row.  It is not a claim that no stronger
arithmetic theorem, supplied with additional source data, can identify or
transport the missing finite information.
\end{proof}

\begin{nonclaim}[Decomposition scope]
\label{nonclaim:apparatus:decomposition-scope}
This section does not prove BSD, the Bloch--Kato conjecture, or ETNC~\citep{BurnsFlach2001}.  It
defines the five-column component map, proves the component-killing
consequence of the bridge-defect equation, and records the finite scalar
public-shadow no-go for the typed bridge.  The no-go is an audit result
about what the scalar row fails to determine; it is not a mathematical
non-existence claim about possible richer finite-source identities.  The
external transports from the five native columns to the determinant-line
target remain recognition sources carried separately, not derived theorems
of this decomposition section.
\end{nonclaim}

\section{Height/regulator column}
\label{sec:height-regulator}

The height/regulator column keeps the Mordell--Weil lattice and its
N\'eron--Tate pairing~\citep{Neron1965,Silverman2009} before any comparison
with the global determinant line.  Its readout is the regulator, viewed as
the coefficient of a determinant.

\begin{definition}[Real height-regulator map]\label{def:apparatus:height-regulator-map}
Let
\[
  \Lambda_E=E(\Q)/E(\Q)_{\tors},
  \qquad
  V_E=\Lambda_E\otimes_{\Z}\R,
\]
and let $\langle\cdot,\cdot\rangle_{\mathrm{NT}}$ be the N\'eron--Tate
pairing on $V_E$, with associated map
\[
  h_{\mathrm{NT}}:V_E\longrightarrow V_E^*,
  \qquad
  P\longmapsto \langle P,\cdot\rangle_{\mathrm{NT}}.
\]
The real height-regulator map sends the triple
$(V_E,\Lambda_E,\langle\cdot,\cdot\rangle_{\mathrm{NT}})$ to
\[
  \rho_{\mathrm{ht/reg}}^{\R}
  =\det(h_{\mathrm{NT}})\in\det(V_E^*)\otimes\det(V_E)^*.
\]
If $(P_1,\ldots,P_r)$ is a $\Z$-basis of $\Lambda_E$ and
$H_E=\bigl(\langle P_i,P_j\rangle_{\mathrm{NT}}\bigr)_{1\le i,j\le r}$ is
the Gram matrix, the coefficient of this element in the induced basis is
$\det H_E=\Reg(E/\Q)$, with the empty determinant equal to $1$ when $r=0$.
The coefficient does not depend on the basis: a change of basis by
$A\in\mathrm{GL}_r(\Z)$ replaces $H_E$ by $A^{\mathsf T}H_EA$, and
$\det(A^{\mathsf T}H_EA)=\det(A)^2\det H_E=\det H_E$ because
$\det A=\pm1$.  (A change of basis of finite index $m>1$ multiplies the
coefficient by $m^2$, so integrality of the basis matters.)
\end{definition}

The adequacy test for this column asks whether the full height matrix is
enough to recover the first-order variation of $\log\Reg$.  It is, for a
simple reason: the readout is a linear function of the variation, and the
full source sees every coordinate of the variation.

\begin{theorem}[Height tangent Schur collapse]\label{thm:apparatus:height-schur-collapse}
Let $r\ge1$ and let $H$ be a positive definite real symmetric $r\times r$
matrix.  Identify a symmetric variation $\delta H$ with its coordinate
vector $x=\operatorname{vech}(\delta H)\in\R^n$, $n=r(r+1)/2$, listing the
entries $\delta h_{ij}$ with $i\le j$.  Then
\[
  D_H(\delta H)
  =
  d(\log\det)_H(\delta H)
  =
  \tr(H^{-1}\delta H)
  =
  d_Hx,
\]
where the row $d_H$ has entry $(H^{-1})_{ii}$ at each diagonal
coordinate and $2(H^{-1})_{ij}$ at each off-diagonal coordinate $i<j$.
Take the audit energy $C=I_n$ on these coordinates and the full source
$L_{\mathrm{full}}(x)=x$.  Then
\[
  \KDD=\|d_H\|^2,
  \qquad
  \XiC=d_Hd_H^{\mathsf T}-d_H\,I_n\,d_H^{\mathsf T}=0 .
\]
For $r=0$ the tangent space is zero, $n=0$, and $\KDD=0$: the rank-zero
row is vacuous.%
\footnote{Lean proves the algebraic core: for every row $d$ of every
length over any commutative ring, the identity source gives Schur
residual zero, and the identity matrix is the unique matrix satisfying
the four Penrose equations at the identity.  For $r=2$ it also proves
that the displayed row computes $\tr(H^{-1}\delta H)$ when
$\det H\neq0$.  The calculus step identifying the row with the derivative
of $\log\det$, and the general-rank row, are proved here on paper.}
\end{theorem}

\begin{proof}
Since $H$ is invertible, $\det(H+t\,\delta H)=\det H\cdot\det(I+tA)$ with
$A=H^{-1}\delta H$.  In the permutation expansion of $\det(I+tA)$, the
identity permutation contributes $1+t\tr A+O(t^2)$, and every other
permutation moves at least two indices and contributes $O(t^2)$.  Since
$\det H>0$, the chain rule gives $d(\log\det)_H(\delta H)=\tr(H^{-1}\delta
H)$.  Writing $G=H^{-1}$, which is symmetric,
\[
  \tr(G\,\delta H)=\sum_{i,j}G_{ij}\,\delta h_{ji}
  =\sum_i G_{ii}\,\delta h_{ii}+\sum_{i<j}2G_{ij}\,\delta h_{ij},
\]
because the off-diagonal coordinate $\delta h_{ij}$ occurs twice in the
symmetric matrix $\delta H$.  This is $d_Hx$.

For the Schur residual, the source block is $\KLL=I_n$, whose
Moore--Penrose inverse is $I_n$, and the cross blocks are $\KDL=d_H$ and
$K_{LD}=d_H^{\mathsf T}$.  Hence
$\KDL\Kdagger K_{LD}=d_H\,I_n\,d_H^{\mathsf T}=d_Hd_H^{\mathsf T}=\KDD$,
and $\XiC=0$.  If $r=0$, the row $d_H$ is empty and the empty sum
$d_Hd_H^{\mathsf T}$ is $0$.
\end{proof}

For example, when $r=2$ and $H=\bigl(\begin{smallmatrix}a&b\\
b&c\end{smallmatrix}\bigr)$ with $\Delta=ac-b^2>0$, the row in the
coordinates $(\delta h_{11},\delta h_{12},\delta h_{22})$ is
\[
  d_H=\Bigl(\frac{c}{\Delta},\,-\frac{2b}{\Delta},\,\frac{a}{\Delta}\Bigr),
\]
as one also sees from $\det(H+t\,\delta H)=\Delta+t\,(c\,\delta h_{11}-2b\,
\delta h_{12}+a\,\delta h_{22})+O(t^2)$.  When $r=1$, $d_H=1/\Reg(E/\Q)$
and $\KDD=\Reg(E/\Q)^{-2}$; for the curve $37a1$, the regulator
$0.0511114\ldots$ recorded in the LMFDB~\citep{LMFDB} gives
$\KDD=382.79\ldots$.  The choice $C=I_n$ refers to the chosen
$\operatorname{vech}$ coordinates; it is not invariant under changes of
lattice basis, and only the vanishing of $\XiC$, not the value of $\KDD$,
is used later.

\begin{obligation}[Height/regulator transport]\label{obl:apparatus:rho-ht-reg-transport}
The Bloch--Kato height/regulator component requires a comparison
\[
  \beta_{\mathrm{BK,ht}}:
  \det(V_E^*)\otimes\det(V_E)^*
  \longrightarrow
  \Delta_f(M_E)_{\mathrm{ht/reg}}
\]
identifying the N\'eron--Tate determinant with the Beilinson--Deligne
regulator normalization~\citep{Beilinson1985}.  It must fix
the coefficient and lattice conventions, the archimedean
Deligne-cohomology normalization, and the orientation and
trivialization, and it must be compatible with the finite and
determinant-assembly columns without absorbing them.  This comparison is
external arithmetic input.  It is not derived from the finite calculation
above, which only shows that the full height matrix is adequate for the
first-order regulator readout.
\end{obligation}

\begin{warning}[Rank-zero carriers]\label{warn:apparatus:571a1-carrier}
For a curve of rank $0$ the height/regulator column is empty: there is no
nonzero height variation, and the regulator is the empty determinant $1$.
The curve $571a1$, which reappears in \Cref{sec:finite-source} because
its Tate--Shafarevich group has nontrivial $2$-primary part, is such a
curve: a certified $2$-descent bounds its rank above and below by $0$
(details in \Cref{sec:finite-source}).  It should not be confused with
the rank-two curve $571b1$ of the same conductor.  Rank-two
height/regulator examples in this paper use $389a1$.
\end{warning}

\section{Analytic/period column}
\label{sec:analytic-period}

The analytic/period column keeps the analytic leading coefficient and the
real period apart from the height, finite and local data.  Its native data
form a pair of numbers, and its readout is their quotient.

\begin{definition}[Real analytic/period map]\label{def:apparatus:analytic-period-map}
By modularity~\citep{Wiles1995,BreuilConradDiamondTaylor2001},
$L(E,s)=L(f_E,s)$ extends to an entire function.  Let
$r_{\mathrm{an}}=\operatorname{ord}_{s=1}L(E,s)$ and
\[
  A_E=\frac{L^{(r_{\mathrm{an}})}(E,1)}{r_{\mathrm{an}}!}\neq0
\]
be the analytic leading coefficient, that is, the leading Taylor
coefficient of $L(E,s)$ at $s=1$.  Let $\Omega_E^+>0$ be the real period
in the BSD normalization: the integral of $|\omega_E|$ over $E(\R)$ for a
N\'eron differential $\omega_E$.  The native analytic/period family is the
pair
\[
  L_{\mathrm{an/per}}(E)=(A_E,\Omega_E^+),
\]
and the real analytic/period map is
\[
  \rho_{\mathrm{an/period}}^{\R}:(A_E,\Omega_E^+)\longmapsto
  (A_E/\Omega_E^+)\,e_{\mathrm{an/per}},
\]
where $e_{\mathrm{an/per}}$ is a basis vector of a real line.  The column
contains no height, torsion, Tamagawa or $\Sha$ data.
\end{definition}

The adequacy test asks whether the pair is enough to recover the first-order
variation of the quotient.  In logarithmic coordinates the absolute value
of the quotient becomes a difference; its sign is not a tangent direction
and is recorded separately, as part of the orientation data of
\Cref{sec:det-assembly}.

\begin{theorem}[Analytic/period tangent Schur collapse]\label{thm:apparatus:analytic-period-schur-collapse}
In the chart
\[
  x=(x_A,x_\Omega)=(\log|A_E|,\log\Omega_E^+)\in\R^2,
\]
take $C=I_2$, the full native source $L_{\mathrm{full}}(x)=x$, and the
readout $D(x)=x_A-x_\Omega=\log|A_E/\Omega_E^+|=dx$ with $d=(1,-1)$.  Then $\KLL=I_2$,
$\KDD=dd^{\mathsf T}=2$ and
\[
  \XiC=2-d\,I_2\,d^{\mathsf T}=0.
\]
With no source admitted the residual is $2$; with only $x_A$ admitted it
is $1$; with both coordinates admitted it is $0$.%
\footnote{Lean proves these three integer residual computations for the
row $(1,-1)$.}
\end{theorem}

\begin{proof}
The readout energy is $\KDD=1^2+(-1)^2=2$.  With the full source,
$\KLL=\Kdagger=I_2$, $\KDL=d$ and $K_{LD}=d^{\mathsf T}$, so the projected
energy is $d\,I_2\,d^{\mathsf T}=2$ and $\XiC=0$.  With no source nothing is
projected away and the residual is $2$.  If only $x_A$ is admitted, $d$
splits orthogonally as the seen part $(1,0)$ plus the unseen part
$(0,-1)$, and the residual is the squared length $1$ of the unseen part.
\end{proof}

The theorem is elementary, and that is the point: a single observed
coordinate cannot recover a quotient of two independent quantities,
while the pair can.  The arithmetic content of the column lies entirely in
the comparison below.

\begin{obligation}[Analytic/period transport]\label{obl:apparatus:rho-an-period-transport}
The Bloch--Kato analytic component requires a comparison
\[
  \beta_{\mathrm{BK,an}}:
  L_{\mathrm{an/per}}(E)_{\R}
  \longrightarrow
  \Delta_f(M_E)_{\mathrm{an/per}}
\]
identifying $A_E/\Omega_E^+$ with the zeta-element factor of the
fundamental line.  It must fix the rational and integral normalizations,
the Betti--de Rham comparison and the real orientation, the leading-term
convention in positive rank, and the separation from the other four
columns.  We treat this comparison as external input in every rank.  In
rank $0$ the analytic quotient $L(E,1)/\Omega_E^+$ is a classical object,
but matching it to the zeta-element factor in the conventions above is
part of the comparison, not a consequence of the Schur calculation.
\end{obligation}

\section{Finite-source column and the finite no-gos}
\label{sec:finite-source}

The finite-source column keeps the arithmetic finite data before it is
contracted to the public scalar finite factor.  Its typed tower separates
Tamagawa factors~\citep{Silverman1994}, torsion, the $p$-primary \(\Sha\) data, and the
Cassels--Tate pairing.  The point of this column is that the determinant
readout is pairing-aware: the Cassels--Tate form~\citep{PoonenStoll1999,Cassels1962} is source data, not a
number to be recovered from the scalar finite product after the fact.

\begin{definition}[Real finite-source map]\label{def:apparatus:finite-source-map}
The finite-source tower is
\[
  C_{\mathrm{fin}}(E)
  =
  C_{\mathrm{Tam}}\oplus C_{\tors}\oplus
  C_{\Sha[p^\infty]}\oplus C_{\mathrm{CT}}.
\]
Its native family is
\[
  L_{\mathrm{finite}}(E)
  =
  \bigl((c_v(E))_{v\mid N},\,E(\Q)_{\tors},\,
  (\Sha(E/\Q)[p^\infty])_p,\,\langle\cdot,\cdot\rangle_{\mathrm{CT}}\bigr).
\]
The real finite-source map is the tensor
\[
  \rho_{\mathrm{finite}}^{\R}
  =
  \rho_{\mathrm{Tam}}^{\R}
  \otimes
  \rho_{\tors}^{\R}
  \otimes
  \bigotimes_p \rho_{\Sha\text{-}\mathrm{CT},p}^{\R}.
\]
The Tamagawa factor is
\[
  \rho_{\mathrm{Tam}}^{\R}:
  (c_v)_{v\mid N}
  \longmapsto
  \bigotimes_{v\mid N} c_v e_v^{\mathrm{loc}},
  \qquad
  c_v=|\Phi_v(E)|,\quad
  \Phi_v=E(\Q_v)/E^0(\Q_v).
\]
The torsion factor is
\[
  \rho_{\tors}^{\R}:
  T(E)
  \longmapsto
  \bigotimes_p |T_p(E)|^{-2}e_p^{\tors}.
\]
For the \(\Sha\)--Cassels--Tate factor, put
\[
  S_p=\Sha[p^\infty]/\Sha_{\mathrm{div}}[p^\infty],
\]
equipped with its perfect alternating Cassels--Tate pairing
\(\langle\cdot,\cdot\rangle_{\mathrm{CT},p}\).  Then
\[
  \rho_{\Sha\text{-}\mathrm{CT},p}^{\R}:
  (S_p,\langle\cdot,\cdot\rangle_{\mathrm{CT},p})
  \longmapsto
  \Pf(\langle\cdot,\cdot\rangle_{\mathrm{CT},p})
  e_p^{\Sha\text{-}\mathrm{CT}}.
\]
For a trivial \(S_p\), the Pfaffian contribution is the unit.  The source
datum in this block is the pair
\((S_p,\langle\cdot,\cdot\rangle_{\mathrm{CT},p})\), not merely
\(|S_p|\) and not merely \(\dim_{\mathbb F_p}S_p[p]\).
\end{definition}

\begin{theorem}[Pairing-aware finite Schur collapse]\label{thm:apparatus:finite-source-schur-collapse}
At a support prime $p$, use finite tangent coordinates
\[
  x=(x_T,x_{\tors},x_{\Sha},x_{\mathrm{CT}})\in\R^4,
\]
where \(x_T=v_p(\prod_v c_v)\), \(x_{\tors}=v_p|E(\Q)_{\tors}|\),
\(x_{\Sha}=v_p|S_p|\), and \(x_{\mathrm{CT}}\) is the
Cassels--Tate-pairing coordinate.  With the pairing-aware full readout
\[
  D_{\mathrm{full}}(x)=x,
  \qquad
  L_{\mathrm{full}}(x)=x,
  \qquad
  C=I_4,
\]
the Schur residual collapses:
\[
  \XiC
  =
  I_4-I_4I_4^{\dagger}I_4
  =
  0.
\]
Thus the full typed finite tower is adequate for the pairing-aware finite
readout.\footnote{Verified in Lean as
\texttt{SixBirdsBSD.Apparatus.FiniteSource.finiteSourceSchurCollapse}; see App~\ref{app:formalization}.}
\end{theorem}

\begin{proof}
Because the full readout and the full source are both the identity on
\(\R^4\), the readout block is
\[
  \KDD=I_4.
\]
The source block is also \(\KLL=I_4\).  Since the Moore--Penrose inverse of
the identity matrix is the identity matrix, \(\Kdagger=I_4\).  The two cross
blocks are \(\KDL=I_4\) and \(K_{LD}=I_4\).

It remains to compute the projection term rather than name it.  For any
vector \(v=(v_1,v_2,v_3,v_4)^{\mathsf T}\),
\[
  I_4v=v.
\]
Equivalently, in coordinates,
\[
  (I_4v)_i=\sum_{j=1}^4\delta_{ij}v_j=v_i.
\]
Applying the same identity-matrix action column by column to \(I_4\) gives
\[
  I_4I_4=I_4.
\]
Therefore
\[
  \KDL\Kdagger K_{LD}
  =
  I_4I_4I_4
  =
  I_4.
\]
Substitution into the Schur-complement residual gives
\[
  \XiC
  =
  \KDD-\KDL\Kdagger K_{LD}
  =
  I_4-I_4
  =
  0.
\]
This calculation uses all four finite coordinates.  The scalar and
dimension-only shadows treated below are different readouts: they forget
part of the typed finite tower and therefore leave positive residuals.
\end{proof}

\begin{theorem}[Scalar finite projection (NG-1, Schur form)]\label{thm:apparatus:scalar-finite-shadow}
Let \(x=(x_T,x_{\tors},x_{\Sha},x_{\mathrm{CT}})\in\R^4\), and let the
scalar finite projection be
\[
  q=s\cdot x,
  \qquad
  s=(1,-2,1,0).
\]
The orthogonal rank-one projector onto the scalar row is
\[
  P_s=s^{\mathsf T}(ss^{\mathsf T})^{-1}s,
\]
and the residual against the full typed finite readout is
\[
  \Xi_{\mathrm{scalar}}=I_4-P_s,
  \qquad
  \tr(\Xi_{\mathrm{scalar}})=3>0.
\]
Moreover the scalar map is noninjective on the typed finite tower: the typed
finite inputs \((0,0,2,H)\) and \((2,0,0,0)\) both give \(q=2\), but place
the valuation in the \(\Sha\)--Cassels--Tate block and the Tamagawa block,
respectively.\footnote{Verified in Lean as
\texttt{SixBirdsBSD.Apparatus.FiniteSource.scalarFiniteShadow}; see App~\ref{app:formalization}.}
\end{theorem}

\begin{proof}
First compute the Schur residual.  The squared length of the scalar row is
\[
  ss^{\mathsf T}
  =
  1^2+(-2)^2+1^2+0^2
  =
  6.
\]
Thus the rank-one projector has entries
\[
  (P_s)_{ij}=\frac{s_i s_j}{6}.
\]
Its trace is the sum of its diagonal entries:
\[
  \tr(P_s)
  =
  \sum_{i=1}^4 \frac{s_i^2}{6}
  =
  \frac{1^2+(-2)^2+1^2+0^2}{6}
  =
  \frac{6}{6}
  =
  1.
\]
Since \(\tr(I_4)=4\), the residual trace is
\[
  \tr(\Xi_{\mathrm{scalar}})
  =
  \tr(I_4-P_s)
  =
  \tr(I_4)-\tr(P_s)
  =
  4-1
  =
  3.
\]
This is positive, so the scalar row leaves a three-dimensional residual
relative to the full four-coordinate finite source.

The same failure is visible without linear algebra.  Evaluate \(q\) on the
two displayed typed inputs.  For the \(\Sha\)--Cassels--Tate input
\((0,0,2,H)\),
\[
  q=1\cdot0-2\cdot0+1\cdot2+0\cdot H=2.
\]
For the Tamagawa input \((2,0,0,0)\),
\[
  q=1\cdot2-2\cdot0+1\cdot0+0\cdot0=2.
\]
The scalar value is the same, but the typed data are different: one input
records the contribution in the \(\Sha\)--Cassels--Tate block, while the
other records it in the Tamagawa block.  The scalar projection also ignores
the Cassels--Tate coordinate \(H\) itself, since the fourth coefficient of
\(s\) is zero.  Hence the scalar finite row is a public shadow of the typed
finite tower.  This is an audit no-go: it says the scalar row alone does not
recover the typed source; it is not a claim that no richer theorem can
supply the missing finite data.
\end{proof}

\begin{theorem}[{$\dim\Sha[p]$ projection (NG-2, Schur form)}]\label{thm:apparatus:dim-sha-shadow}
On the \(\Sha\)--Cassels--Tate subcarrier \(y=(m,\theta)\in\R^2\), where
\(m\) records the \(p\)-power elementary-divisor length and \(\theta\)
records the Cassels--Tate coordinate, the dimension-only projection
\[
  L_{\dim}(y)=(1,0)y
\]
leaves
\[
  \Xi_{\dim}
  =
  I_2-(1,0)^{\mathsf T}(1,0)
  =
  \operatorname{diag}(0,1),
  \qquad
  \tr(\Xi_{\dim})=1>0.
\]
It is noninjective on Cassels--Tate data: \((\Z/2\Z)^2\) with
\(\langle e,f\rangle=\tfrac12\) and \((\Z/4\Z)^2\) with
\(\langle e,f\rangle=\tfrac14\) have the same
\(\mathbb F_2\)-dimension of their \(2\)-torsion but different
Cassels--Tate determinant factors.\footnote{Verified in Lean as
\texttt{SixBirdsBSD.Apparatus.FiniteSource.dimShaShadow}; see App~\ref{app:formalization}.}
\end{theorem}

\begin{proof}
Let \(e=(1,0)\).  The dimension-only source keeps the \(m\)-coordinate and
forgets the \(\theta\)-coordinate.  Its projector is
\[
  e^{\mathsf T}e
  =
  \begin{pmatrix}1&0\\0&0\end{pmatrix}.
\]
Therefore the residual projector is
\[
  \Xi_{\dim}
  =
  I_2-e^{\mathsf T}e
  =
  \begin{pmatrix}1&0\\0&1\end{pmatrix}
  -
  \begin{pmatrix}1&0\\0&0\end{pmatrix}
  =
  \begin{pmatrix}0&0\\0&1\end{pmatrix}
  =
  \operatorname{diag}(0,1).
\]
Taking traces gives
\[
  \tr(\Xi_{\dim})=0+1=1>0.
\]
Thus the dimension-only row leaves exactly the Cassels--Tate direction
unexplained.

The witness shows that this residual is not a bookkeeping artifact.  The
group \(U_1=(\Z/2\Z)^2\) has \(U_1[2]=U_1\), so
\(\dim_{\mathbb F_2}U_1[2]=2\).  The group \(U_2=(\Z/4\Z)^2\) has
\[
  U_2[2]\cong(\Z/2\Z)^2,
\]
so \(\dim_{\mathbb F_2}U_2[2]=2\) as well.  The dimension readout therefore
cannot distinguish \(U_1\) from \(U_2\).

Their Cassels--Tate determinant factors differ.  For a two-generator
alternating pairing with matrix
\[
  \begin{pmatrix}0&a\\-a&0\end{pmatrix},
\]
the Pfaffian is \(a\).  The pairings in the statement have \(a=\tfrac12\)
and \(a=\tfrac14\), respectively, so they give different Pfaffian
contributions to the determinant-line trivialization.  Hence the dimension
\(\dim_{\mathbb F_2}\Sha[2]\) is a noninjective public shadow of the
\(\Sha\)--Cassels--Tate source.  As above, this is an audit no-go about a
specific public readout, not a non-existence claim about all possible
arithmetic refinements.
\end{proof}

\begin{obligation}[Refined finite-source BK transport]\label{obl:apparatus:rho-finite-transport}
The full Bloch--Kato finite component requires a comparison transport
\[
  \beta_{\mathrm{BK,finite}}:
  L_{\mathrm{finite}}(E)_{\R}
  \longrightarrow
  \Delta_f(M_E)_{\mathrm{finite}}.
\]
This transport fixes the integral local-condition conventions at bad primes,
the squared-torsion normalization, the Cassels--Tate Pfaffian and orientation
conventions, and compatibility with the determinant-assembly column.  In
this paper it is a recognition-source carrier: the finite-source transport
is named external arithmetic content and carried separately, not derived
from the finite typed Schur calculation.
\end{obligation}

% ===========================================================================

\section{p-adic column}
\label{sec:p-adic}

The local \(p\)-adic column records the local source data before it is
contracted to a public local scalar.  At an ordinary prime this means keeping
the ordinary-control coordinate together with the finite local condition and
dual exponential.  At a supersingular prime it means keeping the signed
\(+\) and \(-\) local directions separately, rather than replacing them by
their unsigned aggregate.

\begin{definition}[Real p-adic map]\label{def:apparatus:p-adic-map}
Let \(T=T_pE\) and \(V=V_p(E)=T\otimes\Qp\), viewed as a representation of
\(G_{\Qp}\).  The finite local condition~\citep{BlochKato1990,Fontaine1994} is
\[
  \Hf(\Qp,V)
  =
  \ker\bigl(H^1(\Qp,V)\to H^1(\Qp,V\otimes B_{\mathrm{cris}})\bigr),
\]
and the dual exponential is
\[
  \exp^*:\Hf(\Qp,V)\longrightarrow
  D_{\mathrm{dR}}(V)/\operatorname{Fil}^0D_{\mathrm{dR}}(V).
\]

If \(E\) has good ordinary reduction at \(p\)~\citep{Mazur1972}, so \(a_p\) is a \(p\)-unit
and the ordinary filtration is
\[
  0\longrightarrow T^+\longrightarrow T\longrightarrow T^-\longrightarrow0,
\]
then the native ordinary family is
\[
  L_p^{\mathrm{ord}}(E)=(T^+,\Hf(\Qp,V),\exp^*),
\]
and the real ordinary local map is
\[
  \rho_{p,\mathrm{ord}}^{\R}:
  (T^+,\Hf(\Qp,V),\exp^*)
  \longmapsto
  e_{T^+}\otimes\det(\exp^*)\,e_{\mathrm{loc},p}
  \in \Delta_f(V,T)_{p,\R}^{\mathrm{ord}}.
\]
If \(E\) has good supersingular reduction at \(p\)~\citep{Kobayashi2003,Pollack2003}, with \(a_p\equiv0\),
then the signed native family is
\[
  L_p^{\pm}(E)=(H^1_+,H^1_-,\exp^*),
\]
and the signed real local map is
\[
  \rho_{p,\pm}^{\R}:
  (H^1_+,H^1_-,\exp^*)
  \longmapsto
  e_+\otimes e_-\otimes\det(\exp^*)\,e_{\mathrm{loc},p}^{\pm}.
\]
\end{definition}

\begin{theorem}[Ordinary local Schur collapse]\label{thm:apparatus:ordinary-schur-collapse}
For a good-ordinary row, let
\[
  x=(x_{\mathrm{ord}},x_f,x_{\exp})\in\R^3
\]
be the ordinary-control, finite-local-condition, and dual-exponential
tangent coordinates.  With
\[
  D_{\mathrm{ord}}(x)=x,\qquad
  L_{\mathrm{ord,full}}(x)=x,\qquad
  C=I_3,
\]
one has \(\XiC=0\).  If the ordinary-control coordinate
\(x_{\mathrm{ord}}\) is omitted, the trace residual is \(1\); the unit-root
control coordinate is therefore necessary for the full ordinary local
source.\footnote{Verified in Lean as
\texttt{SixBirdsBSD.Apparatus.PAdic.ordinarySchurCollapse}; see App~\ref{app:formalization}.}
\end{theorem}

\begin{proof}
With the full ordinary source admitted, both the readout and source maps are
the identity on \(\R^3\).  Hence
\[
  \KDD=I_3,\qquad
  \KLL=I_3,\qquad
  \Kdagger=I_3,\qquad
  \KDL=I_3,\qquad
  K_{LD}=I_3.
\]
The projection term is computed by the identity-matrix action.  For
\(v=(v_1,v_2,v_3)^{\mathsf T}\),
\[
  (I_3v)_i=\sum_{j=1}^3\delta_{ij}v_j=v_i,
\]
so \(I_3v=v\), and applying this column by column gives \(I_3I_3=I_3\).
Therefore
\[
  \KDL\Kdagger K_{LD}
  =
  I_3I_3I_3
  =
  I_3.
\]
Substituting into the Schur residual gives
\[
  \XiC
  =
  \KDD-\KDL\Kdagger K_{LD}
  =
  I_3-I_3
  =
  0.
\]

Now omit \(x_{\mathrm{ord}}\).  The admitted source is the coordinate
subspace spanned by \(e_f=(0,1,0)\) and \(e_{\exp}=(0,0,1)\).  The
orthogonal projector onto this subspace is
\[
  P_{\mathrm{omit}}
  =
  \operatorname{diag}(0,1,1).
\]
The residual projector is
\[
  I_3-P_{\mathrm{omit}}
  =
  \operatorname{diag}(1,0,0),
\]
whose trace is \(1\).  Thus omitting the ordinary-control coordinate leaves
exactly one unexplained local direction.
\end{proof}

\begin{theorem}[Signed local Schur collapse]\label{thm:apparatus:signed-schur-collapse}
For a good-supersingular row, let
\[
  y=(y_+,y_-,y_{\exp})\in\R^3
\]
be the signed \(+\), signed \(-\), and dual-exponential tangent coordinates.
With
\[
  D_{\pm}(y)=y,\qquad
  L_{\pm,\mathrm{full}}(y)=y,\qquad
  C=I_3,
\]
one has \(\XiC=0\).  The unsigned projection
\[
  L_{\mathrm{uns}}(y)=(y_++y_-,y_{\exp})
\]
has rank \(2\), leaves the difference direction \(y_+-y_-\) invisible, and
has trace residual \(1\).  Thus the signed decomposition is structural;
unsigned aggregation is a public shadow.\footnote{Verified in Lean as
\texttt{SixBirdsBSD.Apparatus.PAdic.signedSchurCollapse}; see App~\ref{app:formalization}.}
\end{theorem}

\begin{proof}
The full signed source is again the identity source on \(\R^3\).  Therefore
the same identity-matrix computation as in the ordinary case gives
\[
  \KDD=I_3,\qquad
  \KLL=I_3,\qquad
  \Kdagger=I_3,\qquad
  \KDL=I_3,\qquad
  K_{LD}=I_3,
\]
and hence
\[
  \XiC
  =
  I_3-I_3I_3I_3
  =
  0.
\]

It remains to compute the unsigned residual.  The unsigned source is the
two-dimensional subspace spanned by
\[
  u_1=(1,1,0),\qquad u_2=(0,0,1).
\]
These two vectors are orthogonal, with \(u_1\cdot u_1=2\),
\(u_2\cdot u_2=1\), and \(u_1\cdot u_2=0\).  Hence the orthogonal projector
onto the unsigned source is
\[
  P_{\mathrm{uns}}
  =
  \frac{u_1u_1^{\mathsf T}}{2}
  +
  u_2u_2^{\mathsf T}.
\]
Its trace is
\[
  \tr(P_{\mathrm{uns}})
  =
  \frac{u_1\cdot u_1}{2}
  +
  u_2\cdot u_2
  =
  \frac{2}{2}+1
  =
  2.
\]
Since \(\tr(I_3)=3\), the unsigned residual has trace
\[
  \tr(I_3-P_{\mathrm{uns}})=3-2=1.
\]

The invisible direction is explicit.  The vector
\[
  w=(1,-1,0)
\]
is orthogonal to \(u_1\) and \(u_2\), so it lies in the residual line.
Equivalently,
\[
  L_{\mathrm{uns}}(1,0,0)=(1,0)=L_{\mathrm{uns}}(0,1,0),
\]
although the two signed inputs are distinct.  Thus unsigned aggregation
forgets the \(+\) versus \(-\) difference, while the signed carrier retains
it.
\end{proof}

\begin{obligation}[Refined p-adic BK-local transport]\label{obl:apparatus:rho-p-transport}
The full Bloch--Kato local component requires a comparison transport
\[
  \beta_{\mathrm{BK},p}:
  \Delta_f(V,T)_{p,\R}
  \longrightarrow
  \Delta_f(M_E)_p.
\]
This transport fixes the integral Fontaine--Laffaille or Wach data, the
local Tate-pairing~\citep{Milne1986} and dual-exponential orientation, the Nekov\'a\v{r}
Selmer-complex local-global compatibility~\citep{Nekovar2006}, the signed normalization at
supersingular primes, and compatibility with the determinant-assembly
column.  In this paper it is a recognition-source carrier: the arithmetic
comparison is named external content and carried separately, not derived
from the finite typed local calculations above.
\end{obligation}

% ===========================================================================

\section{Determinant-assembly column}
\label{sec:det-assembly}

The four arithmetic columns produce four factors.  To obtain an element of
one determinant line one also needs an orientation: the signs, units and
trivialization conventions under which the factors are multiplied.  The
determinant-assembly column holds this orientation as an independent
input.

\begin{definition}[Real determinant-assembly map]\label{def:apparatus:det-assembly-map}
Let
\[
  \lambda_{\mathrm{loc}}(E)
  =
  \rho_{\mathrm{ht/reg}}^{\R}
  \otimes
  \rho_{\mathrm{an/period}}^{\R}
  \otimes
  \rho_{\mathrm{finite}}^{\R}
  \otimes
  \bigotimes_{p\in S(E)}\rho_p^{\R}
\]
be the tensor of the four arithmetic columns over a finite support set
$S(E)$ of primes, with the local factors at primes outside $S(E)$
trivialized as units.  Let $o_{\det}\in C_{\det}(E)$ be an orientation and
trivialization datum, supplied independently of the four columns, and let
$\rho_{\det}^{\R}(o_{\det})\in\R^\times$ be the orientation scalar it
determines (the map on the orientation slot from
\Cref{def:apparatus:bk-component-map}).  The real determinant assembly is
\[
  \mathcal A^{\R}(E)
  =
  \lambda_{\mathrm{loc}}(E)\cdot\rho_{\det}^{\R}(o_{\det})
  \in\Delta_f(M_E)_{\R}.
\]
We call $\mathcal A^{\R}(E)$ the \emph{final assembly}.  It already
contains the four arithmetic factors, and so must not be used again as a
factor alongside them; see \Cref{def:apparatus:five-column-tensor}.
\end{definition}

The adequacy question for this column is whether the four arithmetic
columns, without the orientation, can explain the assembled target.  In
the coordinate model they cannot, and the shortfall is exactly one
direction.

\begin{theorem}[Cascade-completion Schur signature]\label{thm:apparatus:cascade-completion-signature}
Use the tangent coordinates
\[
  x=(x_{\mathrm{ht}},x_{\mathrm{an}},x_{\mathrm{fin}},x_p,x_{\det})\in\R^5 .
\]
Let the four arithmetic columns supply the source row
$\ell=(1,1,1,1,0)$ and let the assembled target have row
$d=(1,1,1,1,1)$.  With $C=I_5$,
\[
  \KDD=dd^{\mathsf T}=5,\qquad
  \KLL=\ell\ell^{\mathsf T}=4,\qquad
  \KDL=d\ell^{\mathsf T}=4,
  \qquad
  \XiC=5-4\cdot4^{-1}\cdot4=1.
\]
The residual is the orientation direction $e_{\det}=(0,0,0,0,1)$.%
\footnote{Lean proves these exact integer values.  The statement concerns
the orientation direction only; it does not assert adequacy of each of the
four arithmetic coordinates.}
\end{theorem}

\begin{proof}
The three blocks are the displayed sums of squares and products of
entries, and the one-dimensional source block $4$ has inverse $4^{-1}$.
Hence $\XiC=5-4=1$.  Since $d=\ell+e_{\det}$ with $\ell\cdot e_{\det}=0$,
projecting $d$ onto $\ell$ leaves $e_{\det}$, of squared length $1$.
\end{proof}

\begin{obligation}[Integral determinant assembly]\label{obl:apparatus:rho-det-assembly}
The full determinant comparison requires, for each orientation datum
$o_{\det}$, a transport
\[
  \beta_{\det\text{-}\mathrm{assembly}}(\,\cdot\,;o_{\det}):
  \Delta_{\mathrm{ht/reg}}
  \otimes
  \Delta_{\mathrm{an/period}}
  \otimes
  \Delta_{\mathrm{finite}}
  \otimes
  \bigotimes_{p\in S(E)}\Delta_p
  \longrightarrow
  \Delta_f(M_E),
\]
together with compatible unit trivializations of the local lines at primes
outside $S(E)$.  It must fix the Knudsen--Mumford determinant
signs~\citep{KnudsenMumford1976}, the local--global compatibility of
Selmer complexes, the convention relating the zeta element to the
motivic trivialization, and compatibility with the four column
transports.  The orientation datum is a parameter of the transport,
distinct from the assembled output.  This transport is external
arithmetic input.  The residual computation above shows only that, in the
coordinate model, the four arithmetic rows do not account for the
orientation direction; it does not show that the orientation varies
independently for actual curves.
\end{obligation}

\section[Comparison with the Bloch--Kato scalar package]{Comparison with the Bloch--Kato scalar package}
\label{sec:comparison}

The preceding sections built the five columns separately.  This section
contracts them back to the usual Bloch--Kato scalar package and asks two
questions: is the typed source adequate for that contraction, and what
does the contraction forget?

\begin{remark}[Relation to the standard presentations]\label{rmk:apparatus:five-column-literature-gap}
The fundamental line and its Tamagawa-number formulation go back to
Bloch--Kato~\citep[\S5]{BlochKato1990} and
Fontaine--Perrin-Riou~\citep{FontainePerrinRiou1994}; the formulation of
the leading value in terms of periods and height determinants, modulo
$\Q^\times$, is surveyed by Nekov\'a\v{r}~\citep{Nekovar1994}.  The
five-slot organization used here is a bookkeeping choice made for the
adequacy tests.  We have not compared it systematically with the
literature and make no claim about whether equivalent decompositions
appear there.
\end{remark}

\begin{definition}[Five-column tensor]\label{def:apparatus:five-column-tensor}
The typed five-column tensor is
\[
  T_{\mathrm{5col}}(E)
  =
  \rho_{\mathrm{ht/reg}}^{\R}(E)
  \otimes
  \rho_{\mathrm{an/period}}^{\R}(E)
  \otimes
  \rho_{\mathrm{finite}}^{\R}(E)
  \otimes
  \bigotimes_{p\in S(E)}\rho_p^{\R}(E)
  \otimes
  \rho_{\det}^{\R}(o_{\det}),
\]
whose fifth factor is the orientation scalar of the independent
orientation datum (\Cref{def:apparatus:det-assembly-map}), not the final
assembly $\mathcal A^{\R}(E)$.  Its source family consists of the five
separated logarithmic tangent coordinates
$(x_{\mathrm{ht}},x_{\mathrm{an}},x_{\mathrm{fin}},x_p,x_{\det})$, where
$x_{\det}=\log|\rho_{\det}^{\R}(o_{\det})|$.  The
\emph{orientation-normalized scalar readout} is
\[
  \log\bigl|T_{\mathrm{5col}}(E)\big/\rho_{\det}^{\R}(o_{\det})\bigr|
  =x_{\mathrm{ht}}+x_{\mathrm{an}}+x_{\mathrm{fin}}+x_p ,
\]
the logarithm of the absolute value of the product of the four arithmetic
factors.  It is distinct from the full contraction
$\log|T_{\mathrm{5col}}(E)|$, whose row is $(1,1,1,1,1)$.
\end{definition}

The distinction in the definition matters.  Since $\mathcal A^{\R}(E)$
already contains the four arithmetic factors, using it as the fifth factor
would count each of them twice.  For instance, with arithmetic factors
$2,3,5,7$ and orientation scalar $1$, the final assembly is $210$; the
tensor with the orientation as fifth factor is again $210$, whereas
reusing the final assembly would give $210^2=44100$.  (Lean checks these
values.)

\begin{theorem}[Five-column to Bloch--Kato scalar contraction]\label{thm:apparatus:five-col-bk-comparison}
In the coordinates
$x=(x_{\mathrm{ht}},x_{\mathrm{an}},x_{\mathrm{fin}},x_p,x_{\det})\in\R^5$,
the orientation-normalized scalar readout of
\Cref{def:apparatus:five-column-tensor} is
$\Phi_{\mathrm{5col}\to\mathrm{BK}}(x)=d\cdot x$ with $d=(1,1,1,1,0)$.
With the full typed source $L(x)=x$ and $C=I_5$,
\[
  \Xi_{\mathrm{5col}\to\mathrm{BK}}
  =
  \|d\|^2-d\,I_5\,d^{\mathsf T}
  =
  4-4=0,
\]
so the typed source is adequate for this readout.  The orientation
direction $e_{\det}=(0,0,0,0,1)$ lies in the kernel of
$\Phi_{\mathrm{5col}\to\mathrm{BK}}$: the readout forgets at least that
coordinate.%
\footnote{Lean proves the contraction row, the zero residual and the
kernel vector.  Identifying this coordinate readout with the arithmetic
Bloch--Kato scalar requires the five transports of the column sections
and, in addition, a convention under which that scalar is independent of
the orientation datum.}
\end{theorem}

\begin{proof}
With the identity source, $\KLL=\Kdagger=I_5$, $\KDL=d$ and
$K_{LD}=d^{\mathsf T}$, so the projected energy is
$d\,I_5\,d^{\mathsf T}=dd^{\mathsf T}=4$, which equals $\KDD$; hence the
residual is $0$.  Finally $d\cdot e_{\det}=0$, so changing only
$x_{\det}$ changes the typed coordinate without changing the scalar
readout.
\end{proof}

The two halves of the theorem fit together.  The orientation-normalized
readout is computed correctly from the typed source, so nothing is lost
in that direction; but the typed source carries an orientation coordinate
that this readout does not see, while the full tensor does (changing the
orientation scalar from $1$ to $2$ in the example above changes the
tensor from $210$ to $420$).  To the extent that the arithmetic
Bloch--Kato scalar is orientation-normalized in this sense, which is an
additional convention not supplied by the transports alone, it is a
projection of the typed atlas.

\section{Higher-rank no-go and the rank-2 Kummer source}
\label{sec:higher-rank-no-go}

The higher-rank no-go records a boundary of the height/regulator column.
Once a Mordell--Weil basis is known, the Neron--Tate Gram matrix computes
the regulator determinant.  It does not, by itself, construct independent
motivic-cycle sources in rank at least \(2\).  This is an audit result about
what the height matrix fails to determine; it is not a non-existence claim
about generalized Gross--Zagier~\citep{GrossZagier1986}, Beilinson--Kato~\citep{Beilinson1985,Kato2004}, \(K_2(E)\), or
higher-Chow constructions~\citep{Bloch1986} supplied by other means.

\begin{definition}[Rank-2 Mordell--Weil/Kummer source]\label{def:apparatus:rank-two-kummer-source}
Let
\[
  \Lambda_E=E(\Q)/E(\Q)_{\tors}.
\]
For a prime \(p\), the Kummer map
\[
  \kappa_p:\Lambda_E\otimes\Qp\longrightarrow \Hf(\Q,V_p(E))
\]
sends Mordell--Weil classes to finite Selmer classes.  For the curve
\(389a1\), fix a Mordell--Weil basis \((P_1,P_2)\).  The rank-\(2\)
Mordell--Weil/Kummer candidate is
\[
  c_2^{\mathrm{MW}}(P_i)=\kappa_p(P_i),\qquad i=1,2,
\]
viewed after tensoring with \(\Qp\).  This is a classical rank-\(2\)
Selmer source once the basis is known.  It is not, by this definition
alone, a generalized Heegner-cycle construction or a motivic-cycle source.
\end{definition}

\begin{theorem}[Rank-2 height Schur collapse]\label{thm:apparatus:rank-two-height-schur-collapse}
For the \(389a1\) rank-\(2\) Gram representative, write the symmetric
height variation in coordinates
\[
  (\delta h_{11},\delta h_{12},\delta h_{22})
\]
and write
\[
  d(\log\det)_H
  =
  \ell
  =
  (3.12679\ldots,-0.76771\ldots,2.14483\ldots).
\]
With the full symmetric height-matrix source, \(\KLL=I_3\) and
\[
  \KDD=\|\ell\|^2=14.96650\ldots.
\]
The Schur residual vanishes:
\[
  \Xi_{NT\to\Reg}
  =
  \|\ell\|^2-\ell I_3\ell^{\mathsf T}
  =
  0.
\]
Thus the full rank-\(2\) Neron--Tate Gram carrier is adequate for its
regulator determinant.  By contrast, the scalar-determinant-only source is
one-dimensional and leaves trace residual \(3-1=2\).\footnote{Verified in
Lean as
\texttt{SixBirdsBSD.Apparatus.HigherRankNoGo.rankTwoHeightSchurCollapse};
see App~\ref{app:formalization}.}
\end{theorem}

\begin{proof}
The tangent space of symmetric \(2\times2\) height matrices has the three
coordinates
\[
  (\delta h_{11},\delta h_{12},\delta h_{22}).
\]
The differential of the regulator determinant in logarithmic form is the
row vector \(\ell\).  Hence the readout energy before admitting a source is
\[
  \KDD
  =
  \ell\ell^{\mathsf T}
  =
  \|\ell\|^2
  =
  14.96650\ldots.
\]
When the full symmetric height-matrix carrier is admitted, the source block
is \(I_3\), its Moore--Penrose inverse is again \(I_3\), and the cross
blocks are \(\KDL=\ell\) and \(K_{LD}=\ell^{\mathsf T}\).  The projection
term is the identity-matrix action:
\[
  \ell I_3\ell^{\mathsf T}
  =
  \ell\ell^{\mathsf T}
  =
  \|\ell\|^2.
\]
Therefore
\[
  \Xi_{NT\to\Reg}
  =
  \KDD-\KDL\Kdagger K_{LD}
  =
  \|\ell\|^2-\|\ell\|^2
  =
  0.
\]

For the scalar-determinant-only comparison, the admitted source is the line
spanned by the nonzero vector \(\ell\) inside the three-dimensional
symmetric tangent space.  The orthogonal projector onto this line is
\[
  P_{\ell}=\frac{\ell^{\mathsf T}\ell}{\ell\ell^{\mathsf T}}.
\]
Its trace is
\[
  \tr(P_{\ell})
  =
  \frac{\ell\ell^{\mathsf T}}{\ell\ell^{\mathsf T}}
  =
  1.
\]
Since \(\tr(I_3)=3\), the scalar-determinant-only residual has trace
\[
  \tr(I_3-P_{\ell})=3-1=2.
\]
Thus the determinant scalar is adequate as a readout of the determinant
itself, but it does not retain the full three-coordinate height matrix.
\end{proof}

\begin{theorem}[Gram matrix is a noninjective shadow (NG-3)]\label{thm:apparatus:gram-matrix-shadow}
For \(r\geq2\), the Neron--Tate Gram matrix is a noninjective public shadow
of the rank-\(r\) motivic-cycle source.  The Gram construction records the
pairings
\[
  H_E=(\langle P_i,P_j\rangle_{\mathrm{NT}})_{1\leq i,j\leq r}
\]
for a known Mordell--Weil basis, and therefore supplies
\(\Reg(E/\Q)=\det H_E\).  It does not construct the independent rank-\(r\)
cycle source: distinct motivic-cycle data, including orientation, motivic
labels, and Abel--Jacobi provenance, can have the same Gram matrix.  Thus
the Mordell--Weil height matrix alone cannot be promoted to a rank-\(r\)
cycle-source carrier.\footnote{Verified in Lean as
\texttt{SixBirdsBSD.Apparatus.HigherRankNoGo.gramMatrixShadow}; see App~\ref{app:formalization}.}
\end{theorem}

\begin{proof}
Fix \(r\geq2\).  The Gram construction is a forgetful map.  Its input, if
one is seeking a motivic-cycle source, contains more than the numerical
height pairings: it must also specify which independent source classes are
being realized, their orientation, and their provenance through the relevant
cycle or Abel--Jacobi construction.  The Gram matrix keeps only the numbers
\(\langle P_i,P_j\rangle_{\mathrm{NT}}\).

This forgetfulness gives an explicit noninjectivity witness.  Take the same
Mordell--Weil basis \((P_1,\ldots,P_r)\) and attach two different motivic
source labels or two different Abel--Jacobi provenance records to the same
ordered basis.  These are distinct as motivic-cycle data, because the
source labels or provenance records differ.  Their Neron--Tate pairings are
identical, since the underlying Mordell--Weil points have not changed.
Therefore both data sets map to the same matrix
\[
  (\langle P_i,P_j\rangle_{\mathrm{NT}})_{i,j}.
\]
Equivalently, any height-orthogonal change of representatives that preserves
all pairings also leaves the Gram matrix unchanged while not supplying the
missing motivic labels or provenance.

Once a Mordell--Weil basis is known, the Gram matrix does compute the
regulator: \(\Reg(E/\Q)=\det H_E\).  That is exactly the content used by the
height/regulator column.  The additional claim needed for a rank-\(r\)
cycle source is different: it must construct or recognize independent
motivic classes whose regulator is being measured.  The Gram matrix contains
no such construction.  The no-go is therefore an audit result about the
naive route from height matrix to motivic source; it does not rule out BDP~\citep{BertoliniDarmonPrasanna2013},
Beilinson--Kato, \(K_2(E)\), or higher-Chow constructions imported under
their own hypotheses.
\end{proof}

The rank-\(\geq2\) generalized-Gross--Zagier gap is handled in the
conditional closure of \citet{TsiokosBSDClosure2026}, as a recognition
source rather than a closed theorem.

% ===========================================================================

\section{Finite windows of curves}
\label{sec:global-audit-no-go}

Numerical checks of BSD-type statements are made on finite lists of
curves.  This section records the elementary but easily forgotten fact
that a statement about all curves cannot be decided from a fixed finite
list.  The list used for examples in this paper is the five-curve window
\[
  W=\{11a1,\ 37a1,\ 121b1,\ 960d1,\ 571a1\}.
\]

\begin{theorem}[A finite window does not decide a universal statement]\label{thm:apparatus:finite-window-no-go}
Let $\mathcal C$ be the set of isomorphism classes of elliptic curves over
$\Q$, and consider assignments $\tau:\mathcal C\to\{\mathrm{true},
\mathrm{false}\}$, thought of as truth values of some property of curves.
Let $U(\tau)$ be the statement ``$\tau(E)=\mathrm{true}$ for every $E$''.
\begin{enumerate}
\item No function of the restriction $\tau|_W$ computes $U(\tau)$ for all
  assignments $\tau$: the assignment $\tau_1\equiv\mathrm{true}$ and the
  assignment $\tau_2$ that agrees with $\tau_1$ except at one curve
  $E_*\notin W$, where it is false, have the same restriction to $W$, but
  $U(\tau_1)$ is true and $U(\tau_2)$ is false.
\item In Schur form, let $\mathcal C(B)$ be the finite set of curves of
  conductor at most $B\ge960$, with $M(B)=\#\mathcal C(B)$ coordinates.
  The window source is the coordinate projection onto the five
  coordinates in $W$, and the residual against the identity readout on
  $\mathcal C(B)$ has trace $M(B)-5$, which tends to infinity with $B$.\footnote{Lean proves both parts for a model in which curves are indexed
by natural numbers and the window is the first five indices: the two
completions agree on the window, differ outside it, are separated by the
universal target, and are identified by every function of the window;
and the residual trace $M-5$ takes every natural-number value.  The model
says nothing about which truth assignments arise from actual properties
of curves.}
\end{enumerate}
\end{theorem}

\begin{proof}
(1) Any function of $\tau|_W$ takes the same value on $\tau_1$ and
$\tau_2$, while $U$ separates them.  (2) The trace of a coordinate
projection is the number of retained coordinates, so
$\tr(I_{M(B)}-P_W)=M(B)-5$.  There are infinitely many isomorphism classes
of elliptic curves over $\Q$ and finitely many of each conductor, so
$M(B)\to\infty$.
\end{proof}

The point of the statement is methodological.  A conclusion about all
curves needs an argument that reaches the curves outside the window, for
instance a theorem covering a family, or an exhaustion argument, or a
passage from density statements to pointwise ones.  The finite window by
itself supplies none of these.

\section{Finite sets of primes}
\label{sec:support-prime-no-go}

The local columns are indexed by primes.  This section records the local
analogue of the previous one: a fixed finite set of primes cannot certify
a statement about all primes.  It then describes the numerical support
table used for the examples, and says precisely what that table does and
does not establish.

\begin{theorem}[A finite set of primes does not decide all-prime coverage]\label{thm:apparatus:finite-prime-cover-no-go}
Let a \emph{coverage record} assign to each prime $p$ a status, covered or
uncovered, and let the all-prime target be the statement that every prime
is covered.  For every finite set $P_0$ of primes there is a prime
$q\notin P_0$, and there are two coverage records that agree at every
prime of $P_0$, such that the first covers every prime and the second
leaves $q$ uncovered.  Consequently no function of the restriction of a
record to $P_0$ decides the all-prime target.%
\footnote{Lean proves the statement for coverage records indexed by natural
numbers, with an explicit prime $q$ dividing $1+\prod_{p\in P_0}p$, and
with primality decided by a finite divisor check.  These are formal
records; no claim is made that they arise from elliptic curves.}
\end{theorem}

\begin{proof}
Let $q$ be a prime factor of $1+\prod_{p\in P_0}p$; then $q\notin P_0$.
Let the first record cover every prime and let the second agree with it
except at $q$.  The two records have the same restriction to $P_0$, so
any function of that restriction gives them the same value, while the
all-prime target holds for the first and fails for the second.
\end{proof}

For computations one works with a finite truncation of the local column.

\begin{definition}[Support-prime truncation residual]\label{def:apparatus:support-prime-truncation-residual}
Fix a curve $E$ and a set $R(E)$ of primes, the \emph{residual rows}: the
primes at which no local comparison is taken as available.  For $B\ge1$
the truncation residual is
\[
  \Xi_B(E)=\#\{p\le B:\ p\text{ prime},\ p\in R(E)\}.
\]
It depends only on $R(E)$ below $B$.  In the examples below $R(E)$ is the
set of primes of bad reduction.  For the five window curves and $B=49$
this gives
\[
  \Xi_{49}(11a1)=\Xi_{49}(37a1)=\Xi_{49}(121b1)=1,\qquad
  \Xi_{49}(960d1)=3,\qquad
  \Xi_{49}(571a1)=0,
\]
the last becoming $1$ once $B\ge571$.  This is a finite diagnostic,
distinct from \Cref{thm:apparatus:finite-prime-cover-no-go}.
\end{definition}

Taking $R(E)$ to be the bad primes treats every good prime as covered by
an ordinary or signed local theory.  For good ordinary primes this is the
setting of \Cref{thm:apparatus:ordinary-schur-collapse}.  For good
supersingular primes, the signed theory of
\Cref{def:apparatus:p-adic-map} requires $p$ odd and $a_p=0$, and the
table below shows that this fails in some small cases.

\begin{remark}[Numerical support table]\label{rmk:apparatus:support-prime-atlas}
For the five window curves and the $15$ primes below $50$, together with
the pair $(571a1,571)$, we have computed a $76$-row table of
$(E,p)$ data from the Cremona models: the number of points of the
reduction and $a_p$ (each obtained by two independent methods), the
reduction type (good, multiplicative or additive), and for good primes
whether the reduction is ordinary or supersingular.  The table is
included with the supporting files of the paper.  The bad primes are
\[
  11a1:\{11\},\quad
  37a1:\{37\},\quad
  121b1:\{11\},\quad
  960d1:\{2,3,5\},\quad
  571a1:\{571\}.
\]
Among the good supersingular rows, five lie outside the basic numerical
scope ($p$ odd and $a_p=0$) of Kobayashi's signed
theory~\citep{Kobayashi2003}: $(11a1,2)$ and $(37a1,2)$ with $a_2=-2$;
$(37a1,3)$ with $a_3=-3$; and $(121b1,2)$ and $(571a1,2)$ with $a_2=0$.
For these rows the cited signed input does not supply a local comparison.
Generalized signed constructions that allow $p=2$ or $a_p\neq0$ exist,
but would have to be specified and checked row by row.  If these five
rows are counted as residual, the truncation residuals below $50$ become
$2,3,2,3,1$ for $11a1,37a1,121b1,960d1,571a1$.  The table certifies
the numerical local data only; it does not certify any local comparison
or main-conjecture input for any row.
\end{remark}

The finiteness statements of this paper are collected in
\Cref{tab:no-go-suite}.
\begin{table}[tbp]
\centering
\footnotesize
\caption{The apparatus no-go suite.  Each item is an audit result about what
the displayed public row does not determine; none is a non-existence theorem
about richer external source data.}
\label{tab:no-go-suite}
\begin{tabularx}{\textwidth}{@{}>{\raggedright\arraybackslash}p{0.15\textwidth}
>{\raggedright\arraybackslash}p{0.27\textwidth}
>{\raggedright\arraybackslash}p{0.36\textwidth}
>{\raggedright\arraybackslash}X@{}}
\toprule
No-go & Form & What the literature or public row lacks & Status \\
\midrule
NG-1 & Finite scalar shadow & The scalar finite factor does not identify
which finite-source block carries the valuation. & mechanized finite
noninjectivity witness. \\
\addlinespace
NG-2 & \(\dim \Sha[p]\) shadow & Dimension-only data forgets the
Cassels--Tate determinant and pairing structure. & mechanized projector
trace residual. \\
\addlinespace
NG-3 & Gram-matrix/regulator shadow & A regulator determinant does not
construct the rank-\(2\) source whose Gram matrix it would measure. &
mechanized typed shadow statement. \\
\addlinespace
NG-4 & Finite-window promotion & A bounded audit window cannot certify an
unbounded operator family. & mechanized finite-window obstruction. \\
\addlinespace
NG-5 & Finite-prime coverage & A fixed finite prime set cannot determine an
all-prime support carrier. & mechanized finite-prime cover obstruction. \\
\addlinespace
Higher-rank GZ & Rank-\(\geq2\) source gap & The theorem-grade
generalized Gross--Zagier identity is not supplied by the audited inputs. &
recorded as an audit-result gap. \\
\addlinespace
Global audit & Operator promotion gap & Local or finite scalar evidence does
not by itself promote to the global determinant-line carrier. & recorded as
scope boundary. \\
\addlinespace
Support-prime atlas & Local support gap & The \(b50\) support table is finite
evidence, not all-prime coverage. & support-only audit table. \\
\bottomrule
\end{tabularx}
\end{table}


\section{A formal source algebra in rank two}
\label{sec:vsrc}

The rank-two discussion of \Cref{sec:higher-rank-no-go} suggests separating
two questions: what algebraic shape a rank-$r$ source should have, and
whether actual cycle classes realize that shape.  This section treats the
first question with the simplest possible device, the exterior algebra on
$r$ formal symbols.  The top product $J_1\cdots J_r$ plays the role of a
determinant: it is nonzero exactly when the symbols are independent.  We
call this algebra the \emph{virtual source algebra} $\vsrc_r$; the word
``virtual'' signals that the symbols are not, by themselves, cycle classes.

\begin{definition}[Virtual source algebra]\label{def:apparatus:vsrc-algebra}
Let $V_r=\Q J_1\oplus\cdots\oplus\Q J_r$ and $\vsrc_r=\Lambda_{\Q}(V_r)$,
the exterior algebra.  In rank $2$ it has ordered basis
$1,\ \Jone,\ \Jtwo,\ \Jtwelve:=\Jone\Jtwo$ and relations
\[
  \Jone^2=0,\qquad \Jtwo^2=0,\qquad
  \Jone\Jtwo=\Jtwelve,\qquad
  \Jtwo\Jone=-\Jtwelve.
\]
The element $\Jtwelve$ is the determinant-shaped witness of independence
of $\Jone$ and $\Jtwo$.%
\footnote{Lean constructs the rank-two multiplication table on integer
coordinates and proves associativity, the unit laws and additive and
scalar bilinearity.  Rational coefficients, other ranks and the dimension
count are treated here on paper.}
\end{definition}

\Cref{fig:vsrc-wedge} displays the rank-two algebra.
\begin{figure}[tbp]
\centering
\begin{tikzpicture}[
  font=\small,
  node/.style={draw, rounded corners=2pt, align=center, minimum width=1.6cm,
    minimum height=0.8cm, fill=black!3},
  arrow/.style={-{Latex[length=2mm]}, thick}
]
\node[node] (one) at (0,2.2) {\(1\)};
\node[node] (j1) at (-2.2,0.7) {\(\Jone\)};
\node[node] (j2) at (2.2,0.7) {\(\Jtwo\)};
\node[node] (j12) at (0,-1.1) {\(\Jtwelve\)};

\draw[arrow] (j1) -- node[above, sloped] {\(\wedge\,\Jtwo\)} (j12);
\draw[arrow] (j2) -- node[above, sloped] {\(\wedge\,\Jone=-\)} (j12);
\draw[arrow] (one) -- node[left] {\(\wedge\,\Jone\)} (j1);
\draw[arrow] (one) -- node[right] {\(\wedge\,\Jtwo\)} (j2);
\draw[loop left, looseness=6, -{Latex[length=2mm]}] (j1) to node[left]
  {\(\Jone^2=0\)} (j1);
\draw[loop right, looseness=6, -{Latex[length=2mm]}] (j2) to node[right]
  {\(\Jtwo^2=0\)} (j2);
\node[align=center, font=\footnotesize] at (0,-2.25)
  {\(\dim_{\Q}\vsrc_2=4\), with \(\Jone\Jtwo=\Jtwelve\) the determinant-shaped witness};
\end{tikzpicture}
\caption{The rank-\(2\) \(\vsrc\) exterior-algebra calibration.}
\label{fig:vsrc-wedge}
\end{figure}


\begin{theorem}[Stage I: consistency]\label{thm:apparatus:vsrc-stage-i-consistency}
$\vsrc_2$ is an associative, unital, graded-commutative $\Q$-algebra.
Every element has a unique normal form
$a_{\varnothing}1+a_1\Jone+a_2\Jtwo+a_{12}\Jtwelve$ with $a_\bullet\in\Q$,
so $\dim_{\Q}\vsrc_2=4$.  In lower rank, $\vsrc_0=\Q$ and
$\vsrc_1=\Q\oplus\Q J_1$ with $J_1^2=0$.%
\footnote{Lean proves the unit laws, the four relations and the
distinctness of the basis symbols for the integer-coordinate model, and
records the numbers $2^0,2^1,2^2$; it does not prove the dimension
statement over $\Q$.}
\end{theorem}

\begin{proof}
Take the $\Q$-vector space with basis indexed by the subsets of
$\{1,2\}$.  Define the product of two basis vectors to be $0$ if the
subsets meet, and otherwise to be the basis vector of the union times the
sign of the permutation that sorts the concatenated index list; extend
bilinearly.  This gives the displayed relations, and the empty subset is a
unit.  Associativity holds on basis vectors because both
parenthesizations vanish when an index repeats, and otherwise both
produce the same union with the sign of the same sorting permutation.
Moving a degree-$a$ monomial past a degree-$b$ monomial costs
$(-1)^{ab}$, which is graded commutativity.  The four subsets give a basis,
so the dimension is $4$; the same construction with $\{1\}$ or
$\varnothing$ gives $\vsrc_1$ and $\vsrc_0$.
\end{proof}

\begin{remark}[No rank-$r$ realization is supplied here]\label{rmk:apparatus:vsrc-no-licensed-descent}
For $r\ge2$, nothing in this paper provides a map
\[
  \vsrc_r(E)\longrightarrow \mathrm{CH}^*(\mathcal X)_{\Q},
\]
or into a Selmer or Mordell--Weil group, that realizes $J_1,\ldots,J_r$ as
independent cycle classes.  The Gram-matrix statement
\Cref{thm:apparatus:gram-matrix-shadow} shows that heights alone do not
supply such provenance.  Constructions of cycle classes with arithmetic
meaning, such as Kato's Euler system~\citep{Kato2004} and generalized
Heegner cycles~\citep{BertoliniDarmonPrasanna2013}, come with their own
hypotheses; we do not use them, and we make no claim about whether they
realize $\vsrc_r$ for a given curve.
\end{remark}

In rank $1$ the situation is classical.  Under the Gross--Zagier
hypotheses~\citep{GrossZagier1986,Kolyvagin1990} for an imaginary
quadratic field $K$, there is a Heegner point $P_K\in E(K)$, and one can
send $J_1$ to $P_K$.

\begin{theorem}[Stage II: rank one]\label{thm:apparatus:vsrc-stage-ii-gz-partial}
Send $J_1\mapsto P_K$ and the formal height $h_{\vsrc}(J_1,J_1)$ to the
N\'eron--Tate height $\langle P_K,P_K\rangle_{\mathrm{NT}}$.
\begin{enumerate}
\item The $1\times1$ Gram determinant of the image source is
  $\langle P_K,P_K\rangle_{\mathrm{NT}}$.  If $P_K$ has infinite order and
  $E(K)$ has rank $1$, this equals $[\Lambda:\Z P_K]^2\Reg(E/K)$, where
  $\Lambda=E(K)/E(K)_{\tors}$ and $[\Lambda:\Z P_K]$ is the index of the
  image of $\Z P_K$.
\item The Gross--Zagier identity relating $L'(E/K,1)$ to
  $\langle P_K,P_K\rangle_{\mathrm{NT}}$ is not determined by these
  symbolic data: two formal data sets can share $P_K$ and
  $\langle P_K,P_K\rangle_{\mathrm{NT}}$ while having different analytic
  quotients.  Relative to the symbolic calculus the analytic identity is an
  additional input, which we record by a nonzero residual
  $\Xi_{\mathrm{GZ}}$.\footnote{Lean checks the symbolic assignments $J_1\mapsto P_K$ and
$h_{\vsrc}(J_1,J_1)\mapsto$ (N\'eron--Tate height), and records the
declared residual value $\Xi_{\mathrm{GZ}}=1$; it constructs neither the
Heegner point nor any height.}
\end{enumerate}
\end{theorem}

\begin{proof}
(1) The Gram determinant of a single vector is its height.  If $Q$
generates $\Lambda$ and $P_K\equiv mQ$ modulo torsion, then
$\langle P_K,P_K\rangle_{\mathrm{NT}}=m^2\langle Q,Q\rangle_{\mathrm{NT}}
=m^2\Reg(E/K)$ with $m=[\Lambda:\Z P_K]$.  (2) The symbolic data contain
neither $L'(E/K,1)$ nor the period, so they cannot constrain the analytic
side; the identity is the content of the Gross--Zagier theorem.
\end{proof}

In rank $2$ the corresponding statement is a one-way translation.  For
$389a1$ the Tamagawa product and the torsion order are $1$; we write
\[
  \Xi_{\mathrm{Sc}}
  =A_E-\Omega_E^+\,\Reg(E/\Q)\,\#\Sha(E/\Q)
\]
for the scalar BSD residual, assuming $\Sha(E/\Q)$ finite, which is not
known for this curve.  A \emph{rank-two realization} consists of a
$\Q$-vector space $W$ of classes (for instance cycle classes modulo
homological equivalence), a linear map $f:V_2\to W$, a symmetric
bilinear height pairing on $W$, and a bridge
$\mathcal B^{(2)}_{\mathrm{an\text{-}ht}}$ identifying $A_E$ with
$\Omega_E^+\#\Sha$ times the Gram determinant of $(f(\Jone),f(\Jtwo))$
for that pairing.  Put $r_2=\Lambda f:\vsrc_2\to\Lambda_{\Q}(W)$; then
$z_{12}=r_2(\Jtwelve)=f(\Jone)\wedge f(\Jtwo)$, and $z_{12}\neq0$ exactly
when $f(\Jone),f(\Jtwo)$ are linearly independent.  Let
$\Xi_{\mathrm{Mot}}$ be a \emph{motivic residual}, which by definition
vanishes exactly when a rank-two realization with $z_{12}\neq0$ is given.
All of these data are supplied; none is constructed here.

\begin{theorem}[Stage III: one-way translation for $389a1$]\label{thm:apparatus:vsrc-stage-iii-translation}
For the fixed pair $(\Xi_{\mathrm{Mot}},\Xi_{\mathrm{Sc}})$ just defined,
suppose given the analytic--height bridge in the form of the implication
$\Xi_{\mathrm{Mot}}=0\Rightarrow\Xi_{\mathrm{Sc}}=0$.  Then a rank-two
realization yields the scalar identity
$A_E=\Omega_E^+\Reg(E/\Q)\#\Sha(E/\Q)$.  The reverse implication
$\Xi_{\mathrm{Sc}}=0\Rightarrow\Xi_{\mathrm{Mot}}=0$ is not provided by
the scalar data; it holds only if supplied as separate content.%
\footnote{The formalization records this as a conditional schema: for a
fixed pair of residuals, the forward implication is projected from the
supplied bridge, and the reverse from supplied external content.  Lean
does not construct a realization or a bridge.}
\end{theorem}

\begin{proof}
The forward statement is the supplied implication.  No reverse construction of a realization or
of an independence certificate from the scalar equality is supplied
here.  (\Cref{thm:apparatus:gram-matrix-shadow} establishes only that
provenance data cannot be recovered from the Gram matrix in its record
model; it does not exclude a separately defined construction.)
\end{proof}

The statement must be read for the fixed pair.  A universal version,
asserting $\Xi_{\mathrm{Mot}}=0\Rightarrow\Xi_{\mathrm{Sc}}=0$ for all
values of the two residuals, would be false, for instance at
$(\Xi_{\mathrm{Mot}},\Xi_{\mathrm{Sc}})=(0,1)$.

\begin{remark}[The source reading carries more data]\label{rmk:apparatus:vsrc-rank-r-stronger}
The scalar leading-coefficient formula records a number.  A rank-$r$
source statement in the sense above also asks for a nonzero independent
wedge $J_1\wedge\cdots\wedge J_r$ realized by actual classes.  In this
sense the source reading carries more data than the scalar equality.  We
do not claim a general theorem that scalar BSD is strictly weaker than any
such source statement; the comparison above is for the fixed pair and the
supplied bridge.
\end{remark}

\section[Canonical kappa-r normalization]{$\kapr$ normalization}
\label{sec:kappa-normalization}

The normalization constant \(\kapr\) is the scalar that makes the cascade
leading-term form use the same arithmetic factor as the standard Strong BSD
leading-coefficient formula.  The point of this section is only to identify
that scalar after the cascade convention \(c_E=\Tam(E)\) is fixed.  It is an
equivalence of two ways of writing the same conjectural identity, not a
derivation of BSD.

\begin{theorem}[Canonical $\kapr$ normalization equivalence]\label{thm:apparatus:kappa-r-normalization-equivalence}
Fix the cascade convention
\[
  c_E:=\Tam(E)=\prod_{\ell} c_{\ell}(E).
\]
Then the cascade leading-term form
\[
  \frac{L^{(r)}(E,1)}{r!}
  =
  \kapr(E)\,\Omega_E\,\Reg_r(E)\,c_E
\]
agrees with the standard Strong BSD leading-coefficient form
\[
  \frac{L^{(r)}(E,1)}{r!}
  =
  \Omega_E\,\Reg_r(E)\,
  \frac{\#\Sha(E/\Q)\,\Tam(E)}{\#E(\Q)_{\tors}^{\,2}}
\]
if and only if
\[
  \kapr(E)=
  \frac{\#\Sha(E/\Q)}{\#E(\Q)_{\tors}^{\,2}}.
\]
Under this canonical normalization, the cascade form is equivalent to the
standard Strong BSD identity as a restatement in cascade target coordinates.
It does not prove the identity.\footnote{\raggedright Verified in Lean as
\texttt{SixBirdsBSD.\allowbreak Apparatus.\allowbreak KappaNormalization.\allowbreak kappaRNormalizationEquivalence};
see App~\ref{app:formalization}.}
\end{theorem}

\begin{proof}
After the convention \(c_E=\Tam(E)\), the right-hand side of the cascade
form is
\[
  \kapr(E)\,\Omega_E\,\Reg_r(E)\,\Tam(E).
\]
The right-hand side of the standard Strong BSD leading-coefficient formula
is
\[
  \Omega_E\,\Reg_r(E)\,
  \frac{\#\Sha(E/\Q)\,\Tam(E)}{\#E(\Q)_{\tors}^{\,2}}
  =
  \frac{\#\Sha(E/\Q)}{\#E(\Q)_{\tors}^{\,2}}\,
  \Omega_E\,\Reg_r(E)\,\Tam(E).
\]
Thus the two formulas have the same left-hand side and differ only in the
coefficient multiplying the common factor
\[
  \Omega_E\,\Reg_r(E)\,\Tam(E).
\]
In the BSD normalization this factor is nonzero: the period is positive,
the regulator determinant is positive in positive rank and equal to \(1\)
in rank \(0\), and the Tamagawa product is a positive integer.  Therefore
the two right-hand sides are equal if and only if their scalar coefficients
are equal, namely if and only if
\[
  \kapr(E)
  =
  \frac{\#\Sha(E/\Q)}{\#E(\Q)_{\tors}^{\,2}}.
\]

If this equality of coefficients is imposed, substituting it into the
cascade form gives exactly the standard Strong BSD leading-coefficient
formula.  Conversely, if the cascade form and the standard form agree under
the fixed convention \(c_E=\Tam(E)\), cancellation of the same nonzero common
factor forces the displayed value of \(\kapr(E)\).  The proof is only this
algebraic comparison of normalizations; it supplies no proof that either
leading-coefficient identity is true for \(E\).
\end{proof}

The conditional Strong-BSD closure of \citet{TsiokosBSDClosure2026}
consumes this normalization.  The role of the present section is only to
supply the normalization equivalence used there.

\section{Scope and non-claims}\label{sec:scope}
% statements-of-record rows resolved here: none

The preceding sections build a typed apparatus for comparing the usual
BSD/Bloch--Kato scalar packages with finer source carriers.  The apparatus
is deliberately limited.  It separates what is proved at the finite typed
level from what is supplied by external arithmetic input, and it records
several audit no-gos without turning them into mathematical non-existence
claims.

The four non-claims governing this paper are the following.

\begin{itemize}
  \item \textbf{A-NC-1 --- No-gos are audit results, not non-existence claims.}
  The no-go suite (NG-1--NG-5, higher-rank, global, support-prime)
  records where theorem-grade Mode-A literature does not reach; it is not
  a claim that the missing identities cannot exist.  Here ``Mode-A''
  denotes the theorem-grade arithmetic literature audited as external
  provenance, as opposed to a finite diagnostic, numerical example, or
  support table.

  \item \textbf{A-NC-2 --- Finite typed level.}
  Adequacy collapses are proved at the finite, mathlib-free, typed level
  (identity-matrix action; \(\Q\) rank-1 projector traces);
  operator/motivic full generality is not claimed.  Thus the Schur
  residual calculations in this paper are finite vector and matrix
  computations over explicit typed carriers.  They are not asserted as
  analytic operator theorems or as constructions of motivic cycles.

  \item \textbf{A-NC-3 --- \(\vsrc\) structural narrowing.}
  The \(\vsrc\) exterior-product relations are represented structurally on
  the \(2^r\) coordinate space; the dimension/calibration content is
  faithful, full axiomatization is not claimed.  The \(\vsrc\) section
  therefore calibrates the determinant
  shape of a rank-\(r\) source and the low-rank relations used in the
  paper; it does not present a complete general exterior-algebra
  formalization or a descent to actual motivic cycles.

  \item \textbf{A-NC-4 --- Recognition-source transports are carriers.}
  The five per-column transports are recognition sources (carriers), not
  derived theorems.  They are supplied or imported as named arithmetic
  comparison content and carried as hypotheses where needed; they are not
  established from the finite typed apparatus itself.
\end{itemize}

Consequently, this paper should not be read as proving BSD, an Iwasawa main
conjecture, or a Gross--Zagier theorem.  The column computations prove
adequacy or residual statements for the typed carriers introduced here.
The recognition sources and imports are supplied, imported, or carried as
hypotheses; they are never claimed to be established by the apparatus.

\section{Discussion}\label{sec:discussion}
% statements-of-record rows resolved here: none

The standard BSD formula is a product of scalar factors: an analytic
leading coefficient divided by a period, a regulator, and a finite factor
built from Tamagawa numbers, torsion and $\Sha$, with local $p$-adic
normalizations and determinant conventions hidden in the assembly.  The
formula is a powerful summary, but each factor is the image of richer
data.  The question of this paper is which of those data a given readout
retains.

The Schur residual $\XiC$ turns that question into finite linear algebra.
If all native coordinates of a column are admitted, the residual vanishes.
In general, for the declared target, the residual measures the variation
of the target that the admitted source coordinates do not capture; a
source that discards some coordinates can still be adequate for a target
that ignores them, as the orientation-normalized readout of
\Cref{thm:apparatus:five-col-bk-comparison} shows.  Most of the adequacy
collapses in this paper are elementary for exactly this reason.  Their
value is not depth but bookkeeping: they say which coordinates must be
present before a datum can serve as a source for a given readout.

The ``does not determine'' results complement the collapses.  Some are
formal (\Cref{thm:apparatus:finite-scalar-public-shadow,thm:apparatus:gram-matrix-shadow,thm:apparatus:finite-window-no-go,thm:apparatus:finite-prime-cover-no-go}),
and their content is the precise identification of what is forgotten.
One is arithmetic: the curves $571a1$ and $2045b1$
(\Cref{ex:apparatus:sha-pair}) have the same rank, the same Tamagawa
numbers, the same torsion and the same $2$-torsion dimension of $\Sha$,
but $2$-primary parts of $\Sha$ of orders $4$ and $16$.  The dimension of
$\Sha[2]$ is the datum most directly computed by $2$-descent, and the
example shows concretely that it is not the datum the BSD formula needs.

Three normalization points recur and deserve emphasis.  First, a
Cassels--Tate pairing with values in $\Q/\Z$ does not give a canonical
numerical Pfaffian.  The invariant finite factor is the order of the
paired group; an alternating integral presentation $A$ gives
$\#S_p=|\det A|=\Pf(A)^2$, which determines $|\Pf(A)|$ but not its sign
(\Cref{sec:finite-source}).  Second, on the finite
local condition the relevant map is the logarithm; the dual exponential
vanishes there (\Cref{sec:p-adic}).  Third, the orientation enters the
determinant assembly as an independent input; the assembled value must
not be reused as a factor (\Cref{sec:det-assembly,sec:comparison}).  Each
of these is a small point, but each changes a formula.

At the two ends of the apparatus sit two calibrations.  At the source end,
the exterior algebra $\vsrc_r$ records the shape that a rank-$r$ source
must have, and the rank-two discussion shows that a scalar leading-term
identity does not by itself provide such a source.  At the target end,
the normalization $\kapr=\#\Sha/\#E(\Q)_{\tors}^{\,2}$ is exactly what
makes the cascade form of the leading term coincide with the standard
one.

The language of typed columns, readouts and residuals comes from the Six
Birds framework~\citep{TsiokosNeedles2026,TsiokosFoundationsII2026,TsiokosFoundationsIII2026},
but nothing in the paper depends on that framework beyond the definitions
given here.  The apparatus is used by the companion
paper~\citep{TsiokosBSDClosure2026}, which treats the scalar BSD identity
as a closure problem under explicit hypotheses.

\section{Conclusion}\label{sec:conclusion}
% statements-of-record rows resolved here: none

We separated the data behind the BSD formula into five columns and tested
each column with a finite Schur-complement residual.  The full native data
of each column are adequate for its readout, and the five columns together
are adequate for the orientation-normalized Bloch--Kato scalar, which
forgets the orientation coordinate.  Coarser readouts lose information in
precisely identified ways, and for the $2$-torsion dimension of $\Sha$
this loss occurs for actual curves.  The comparisons from the native
columns to the Bloch--Kato line remain external inputs.  The resulting
picture is a checklist: for each factor of the BSD formula, which source
data must be kept, and which normalizations must be fixed, before the
factor can be compared across theories.

\appendix
\crefalias{section}{appendix}
\section{Lean formalization}\label{app:formalization}
\subsection*{Trust base and methodology}

The Lean development for this paper is mathlib-free.  It imports no
external mathematics library, so no mathlib axioms enter the dependency
closure.  Every declaration registered as a theorem was audited with
\texttt{\#print axioms}.  The only axioms that occur are Lean's standard
foundational ones: \texttt{propext}, \texttt{Quot.sound}, and
\texttt{Classical.choice}, with \texttt{funext} and the \texttt{choice}
recursor appearing as their standard consequences in the audited closures.
There are no project-local axioms behind any theorem-marked result.  Thus
the trust base for the formalized theorem claims is precisely Lean's kernel
plus classical logic, and nothing project-specific.

This limited trust base is possible because the formalized content is
finite and typed.  The adequacy collapses and no-go residuals are not
mechanized as analytic estimates, operator-theoretic assertions, or
motivic constructions.  They are finite vector and matrix computations:
the source and target covariance blocks are explicit finite matrices, the
full-source collapses reduce to identity-matrix action, and the scalar
shadow residuals are computed from rank-\(1\) projectors over \(\Q\).
Each Schur-complement residual is therefore a decidable calculation in a
small finite algebraic model.  This is the reason a mathlib-free kernel can
check the main adequacy and no-go statements.

The surrounding arithmetic geometry is treated separately from those finite
calculations.  The Lean development checks the typed identities appearing
in the apparatus: for example, that an identity source kills the Schur
residual, that the relevant projector traces have the stated values, and
that the finite noninjectivity witnesses have the claimed public shadows.
It does not certify Gross--Zagier, Iwasawa main conjecture inputs,
Cassels--Tate comparison theory, Bloch--Kato local comparison maps, or the
motivic descents named in the body.  Such arithmetic content enters the
paper either as cited external input, as a recognition-source carrier, or
as an explicitly stated hypothesis.

The statements-of-record inventory for this paper contains \(41\) items.
Of these, \(19\) are theorem-grade Lean results with substantive finite
proofs, \(11\) are typed definition mirrors, and \(11\) are not mechanized:
the five recognition-source transport obligations, the four support-only
remarks, the warning, and the scope nonclaim.  Those non-mechanized items
are paper-level statements carried as motivation, scope, or hypotheses,
not Lean declarations.  Two mechanized items, the \(\vsrc\) algebra
definition and the Stage-I \(\vsrc\) consistency theorem, are checked over
a narrowed structural representation; the following subsections record the
per-label coverage, declaration map, and standing narrowings in detail.
\subsection*{Wording-discipline table}

Each theorem in the body carries exactly one mechanization footnote, and
the wording of that footnote is canonical rather than editorial.  The
wording records the fidelity class of the formal counterpart: a faithful
Lean theorem, a Lean theorem over a narrowed representation, a typed
definition mirror, or an unmechanized paper-level statement.  Narrowing and
non-mechanization are therefore visible in the disclosure convention rather
than hidden in prose.

\begin{center}
\begin{tabularx}{\textwidth}{@{}>{\raggedright\arraybackslash}p{0.31\textwidth}X@{}}
\toprule
Statement class & Footnote convention \\
\midrule
Theorem, faithful &
Verified in Lean as \texttt{[decl]}; see App~\ref{app:formalization}. \\
\addlinespace
Theorem, narrowed representation &
Verified in Lean as \texttt{[decl]} over a narrowed representation
(e.g.\ the \(2^r\)-coordinate space); the narrowing is recorded below.  The
narrowing qualifier is mandatory; never bare ``verified in Lean''. \\
\addlinespace
Definition, narrowed representation &
Realized in Lean as \texttt{[decl]} over a narrowed representation; the
narrowing is recorded below. \\
\addlinespace
Definition, faithful (typed mirror) &
No footnote; the declaration appears only in the per-label coverage table
below. \\
\addlinespace
Obligation / remark / warning / non-claim (not mechanized) &
No footnote and no Lean identifier in the body; recorded as a paper-level
statement only. \\
\bottomrule
\end{tabularx}
\end{center}
\subsection*{Standing narrowings}

There are two standing qualifications behind the coverage claims in this
appendix.  Both are part of the formalization contract: they delimit what
the finite Lean development checks, and they keep the surrounding
arithmetic content visible rather than implicit.

First, two \(\vsrc\) items are checked over a narrowed structural
representation: the \(\vsrc\) algebra definition and the Stage-I
\(\vsrc\) consistency theorem.  On the \(2^r\)-coordinate model, the
determinant-shaped content of a rank-\(r\) source is represented, and the
low-rank exterior relations actually used in the paper are checked.  In
rank \(2\), this includes the wedge relation
\[
  \Jone\Jtwo=\Jtwelve
\]
together with the accompanying nilpotence and anticommutation relations.
The dimension count and the calibration of the determinant shape are
faithful to the intended exterior-algebra source.

What is not claimed is equally important.  The formalization is not a
complete general exterior-algebra development: it does not provide general
\(\bigwedge^k\) functoriality, nor a full axiomatization of exterior
algebras in arbitrary rank.  It is also not a descent from \(\vsrc\) to
actual motivic cycles.  The Lean check certifies the finite structural
relations on the coordinate representation, not a general motivic
construction.  This is why those two body items carry the ``narrowed
representation'' footnote convention rather than a bare ``verified in
Lean'' disclosure, and why the \(\vsrc\) statements in
\Cref{sec:vsrc} are calibration and consistency results rather than
constructions of motivic sources.

Second, the five per-column transports are recognition sources, not
derived theorems and not Lean declarations.  These are the
height/regulator, analytic/period, finite-source, \(p\)-adic, and
determinant-assembly comparison transports appearing in
\Cref{sec:height-regulator,sec:analytic-period,sec:finite-source,sec:p-adic,sec:det-assembly}.
Each column's Schur-adequacy collapse is proved at the finite typed level
once the column's native source is admitted.  The corresponding transport
is precisely the named arithmetic comparison content that admits that
source into the Bloch--Kato determinant-line target.  It is supplied or
imported externally and then carried as a hypothesis.

Thus the apparatus does not prove these five transports.  It records them
as obligations and proves the per-column collapses conditional on the
native source data that those transports are meant to recognize.  In the
per-label coverage table below they appear as not-mechanized obligations.
Neither standing narrowing enlarges the Lean trust base: no new axioms are
introduced; the qualifications bound the finite mechanization claims, with
the remaining arithmetic content carried explicitly.
\begin{table}[tbp]
\centering
\small
\caption{Standing representation notes for the apparatus formalization.}
\label{tab:representation-notes}
\begin{tabularx}{\textwidth}{@{}>{\raggedright\arraybackslash}p{0.24\textwidth}
>{\raggedright\arraybackslash}p{0.28\textwidth}
>{\raggedright\arraybackslash}X@{}}
\toprule
Representation issue & Where it appears & Disclosure effect \\
\midrule
\(\vsrc\) structural encoding &
\(\Cref{sec:vsrc}\), App~\ref{app:formalization} &
The rank-\(2\) exterior relations are checked on the \(2^r\)-coordinate
model.  The paper claims calibration and consistency, not a full exterior
algebra or motivic descent formalization. \\
\addlinespace
\(\Q\) rank-\(1\) projector traces &
Finite-source and no-go statements &
Scalar-shadow residuals are finite rational projector computations; the
Lean result is the finite trace calculation, not an arithmetic
classification theorem. \\
\addlinespace
Identity-matrix action &
Schur-collapse statements &
Full-source adequacy reduces to \(I\)-action on an explicit finite carrier.
The external transport that admits the native source remains a carried
recognition obligation. \\
\bottomrule
\end{tabularx}
\end{table}

\subsection*{Per-label coverage}

The following table lists every statement of record in document order,
with its environment kind, its mechanization fidelity, and its body
location.  The paper labels in the first column elide the
\texttt{apparatus:} prefix.  Faithful and narrowed rows have Lean
declarations, mapped in the next subsection; not-mechanized rows do not.

\begin{small}
\begin{longtable}{@{}>{\ttfamily}p{0.45\textwidth} l l l@{}}
\caption{Lean coverage of the labelled statements in the inventory, in document order (41 rows).  Labels omit the \texttt{apparatus:} infix.}\label{tab:formalization-coverage}\\
\toprule \normalfont Label & Kind & Lean coverage & Section \\ \midrule \endfirsthead
\multicolumn{4}{@{}l}{\footnotesize\itshape continued from previous page}\\ \toprule \normalfont Label & Kind & Lean coverage & Section \\ \midrule \endhead
\midrule \multicolumn{4}{r@{}}{\footnotesize\itshape continued on next page}\\ \endfoot
\bottomrule \endlastfoot
def:bk-component-map & \normalfont definition & \normalfont Lean definition & \normalfont \Cref{sec:decomposition} \\
def:bridge-defect-equation & \normalfont definition & \normalfont Lean definition & \normalfont \Cref{sec:decomposition} \\
thm:component-killing & \normalfont theorem & \normalfont Lean proves & \normalfont \Cref{sec:decomposition} \\
thm:finite-scalar-public-shadow & \normalfont theorem & \normalfont model version & \normalfont \Cref{sec:decomposition} \\
nonclaim:decomposition-scope & \normalfont nonclaim & \normalfont not formalized & \normalfont \Cref{sec:decomposition} \\
def:height-regulator-map & \normalfont definition & \normalfont Lean definition & \normalfont \Cref{sec:height-regulator} \\
thm:height-schur-collapse & \normalfont theorem & \normalfont model version & \normalfont \Cref{sec:height-regulator} \\
obl:rho-ht-reg-transport & \normalfont obligation & \normalfont not formalized & \normalfont \Cref{sec:height-regulator} \\
warn:571a1-carrier & \normalfont warning & \normalfont not formalized & \normalfont \Cref{sec:height-regulator} \\
def:analytic-period-map & \normalfont definition & \normalfont Lean definition & \normalfont \Cref{sec:analytic-period} \\
thm:analytic-period-schur-collapse & \normalfont theorem & \normalfont Lean proves & \normalfont \Cref{sec:analytic-period} \\
obl:rho-an-period-transport & \normalfont obligation & \normalfont not formalized & \normalfont \Cref{sec:analytic-period} \\
def:finite-source-map & \normalfont definition & \normalfont Lean definition & \normalfont \Cref{sec:finite-source} \\
thm:finite-source-schur-collapse & \normalfont theorem & \normalfont Lean proves & \normalfont \Cref{sec:finite-source} \\
thm:scalar-finite-shadow & \normalfont theorem & \normalfont Lean proves & \normalfont \Cref{sec:finite-source} \\
thm:dim-sha-shadow & \normalfont theorem & \normalfont model version & \normalfont \Cref{sec:finite-source} \\
obl:rho-finite-transport & \normalfont obligation & \normalfont not formalized & \normalfont \Cref{sec:finite-source} \\
def:p-adic-map & \normalfont definition & \normalfont Lean definition & \normalfont \Cref{sec:p-adic} \\
thm:ordinary-schur-collapse & \normalfont theorem & \normalfont model version & \normalfont \Cref{sec:p-adic} \\
thm:signed-schur-collapse & \normalfont theorem & \normalfont model version & \normalfont \Cref{sec:p-adic} \\
obl:rho-p-transport & \normalfont obligation & \normalfont not formalized & \normalfont \Cref{sec:p-adic} \\
def:det-assembly-map & \normalfont definition & \normalfont Lean definition & \normalfont \Cref{sec:det-assembly} \\
thm:cascade-completion-signature & \normalfont theorem & \normalfont Lean proves & \normalfont \Cref{sec:det-assembly} \\
obl:rho-det-assembly & \normalfont obligation & \normalfont not formalized & \normalfont \Cref{sec:det-assembly} \\
rmk:five-column-literature-gap & \normalfont remark & \normalfont not formalized & \normalfont \Cref{sec:comparison} \\
def:five-column-tensor & \normalfont definition & \normalfont Lean definition & \normalfont \Cref{sec:comparison} \\
thm:five-col-bk-comparison & \normalfont theorem & \normalfont Lean proves & \normalfont \Cref{sec:comparison} \\
def:rank-two-kummer-source & \normalfont definition & \normalfont Lean definition & \normalfont \Cref{sec:higher-rank-no-go} \\
thm:rank-two-height-schur-collapse & \normalfont theorem & \normalfont model version & \normalfont \Cref{sec:higher-rank-no-go} \\
thm:gram-matrix-shadow & \normalfont theorem & \normalfont model version & \normalfont \Cref{sec:higher-rank-no-go} \\
thm:finite-window-no-go & \normalfont theorem & \normalfont model version & \normalfont \Cref{sec:global-audit-no-go} \\
thm:finite-prime-cover-no-go & \normalfont theorem & \normalfont Lean proves & \normalfont \Cref{sec:support-prime-no-go} \\
def:support-prime-truncation-residual & \normalfont definition & \normalfont Lean definition & \normalfont \Cref{sec:support-prime-no-go} \\
rmk:support-prime-atlas & \normalfont remark & \normalfont not formalized & \normalfont \Cref{sec:support-prime-no-go} \\
def:vsrc-algebra & \normalfont definition & \normalfont Lean definition & \normalfont \Cref{sec:vsrc} \\
thm:vsrc-stage-i-consistency & \normalfont theorem & \normalfont model version & \normalfont \Cref{sec:vsrc} \\
rmk:vsrc-no-licensed-descent & \normalfont remark & \normalfont not formalized & \normalfont \Cref{sec:vsrc} \\
thm:vsrc-stage-ii-gz-partial & \normalfont theorem & \normalfont model version & \normalfont \Cref{sec:vsrc} \\
thm:vsrc-stage-iii-translation & \normalfont theorem & \normalfont conditional schema & \normalfont \Cref{sec:vsrc} \\
rmk:vsrc-rank-r-stronger & \normalfont remark & \normalfont not formalized & \normalfont \Cref{sec:vsrc} \\
thm:kappa-r-normalization-equivalence & \normalfont theorem & \normalfont model version & \normalfont \Cref{sec:kappa-normalization} \\
\end{longtable}
\end{small}

\subsection*{Lean declaration map}

The \(30\) mechanized statements, consisting of the \(11\) typed
definitions and \(19\) theorems, each correspond to a Lean declaration
listed here for citation and auditing.  Paper labels elide the
\texttt{apparatus:} prefix, and declarations elide the
\texttt{SixBirdsBSD.Apparatus.} namespace.  The not-mechanized
obligations, remarks, warning, and nonclaim have no declaration and are
therefore omitted; the full coverage table, including those rows, is
\Cref{tab:formalization-coverage}.

\begin{small}
\begin{longtable}{@{}>{\ttfamily\raggedright\arraybackslash}p{0.44\textwidth} >{\ttfamily\raggedright\arraybackslash}p{0.50\textwidth}@{}}
\caption{Lean declaration map for the 30 mechanized apparatus statements (definitions and theorems). Paper labels elide the \texttt{apparatus:} prefix; declarations elide the \texttt{SixBirdsBSD.Apparatus.} namespace.}\label{tab:appd-declmap}\\
\toprule \normalfont Paper label & \normalfont Lean declaration \\ \midrule
\endfirsthead
\multicolumn{2}{@{}l}{\footnotesize\itshape continued from previous page}\\
\toprule \normalfont Paper label & \normalfont Lean declaration \\ \midrule
\endhead
\midrule \multicolumn{2}{r@{}}{\footnotesize\itshape continued on next page}\\
\endfoot
\bottomrule
\endlastfoot
def:\allowbreak bk-\allowbreak component-\allowbreak map & Decomposition.\allowbreak bk\allowbreak Component\allowbreak Map \\
def:\allowbreak bridge-\allowbreak defect-\allowbreak equation & Decomposition.\allowbreak bridge\allowbreak Defect\allowbreak Equation \\
thm:\allowbreak component-\allowbreak killing & Decomposition.\allowbreak component\allowbreak Killing \\
thm:\allowbreak finite-\allowbreak scalar-\allowbreak public-\allowbreak shadow & Decomposition.\allowbreak finite\allowbreak Scalar\allowbreak Public\allowbreak Shadow \\
def:\allowbreak height-\allowbreak regulator-\allowbreak map & Height\allowbreak Regulator.\allowbreak height\allowbreak Regulator\allowbreak Map \\
thm:\allowbreak height-\allowbreak schur-\allowbreak collapse & Height\allowbreak Regulator.\allowbreak height\allowbreak Schur\allowbreak Collapse \\
def:\allowbreak analytic-\allowbreak period-\allowbreak map & Analytic\allowbreak Period.\allowbreak analytic\allowbreak Period\allowbreak Map \\
thm:\allowbreak analytic-\allowbreak period-\allowbreak schur-\allowbreak collapse & Analytic\allowbreak Period.\allowbreak analytic\allowbreak Period\allowbreak Schur\allowbreak Collapse \\
def:\allowbreak finite-\allowbreak source-\allowbreak map & Finite\allowbreak Source.\allowbreak finite\allowbreak Source\allowbreak Map \\
thm:\allowbreak finite-\allowbreak source-\allowbreak schur-\allowbreak collapse & Finite\allowbreak Source.\allowbreak finite\allowbreak Source\allowbreak Schur\allowbreak Collapse \\
thm:\allowbreak scalar-\allowbreak finite-\allowbreak shadow & Finite\allowbreak Source.\allowbreak scalar\allowbreak Finite\allowbreak Shadow \\
thm:\allowbreak dim-\allowbreak sha-\allowbreak shadow & Finite\allowbreak Source.\allowbreak dim\allowbreak Sha\allowbreak Shadow \\
def:\allowbreak p-\allowbreak adic-\allowbreak map & PAdic.\allowbreak p\allowbreak Adic\allowbreak Map \\
thm:\allowbreak ordinary-\allowbreak schur-\allowbreak collapse & PAdic.\allowbreak ordinary\allowbreak Schur\allowbreak Collapse \\
thm:\allowbreak signed-\allowbreak schur-\allowbreak collapse & PAdic.\allowbreak signed\allowbreak Schur\allowbreak Collapse \\
def:\allowbreak det-\allowbreak assembly-\allowbreak map & Det\allowbreak Assembly.\allowbreak det\allowbreak Assembly\allowbreak Map \\
thm:\allowbreak cascade-\allowbreak completion-\allowbreak signature & Det\allowbreak Assembly.\allowbreak cascade\allowbreak Completion\allowbreak Signature \\
def:\allowbreak five-\allowbreak column-\allowbreak tensor & Comparison.\allowbreak five\allowbreak Column\allowbreak Tensor \\
thm:\allowbreak five-\allowbreak col-\allowbreak bk-\allowbreak comparison & Comparison.\allowbreak five\allowbreak Col\allowbreak Bk\allowbreak Comparison \\
def:\allowbreak rank-\allowbreak two-\allowbreak kummer-\allowbreak source & Higher\allowbreak Rank\allowbreak No\allowbreak Go.\allowbreak rank\allowbreak Two\allowbreak Kummer\allowbreak Source \\
thm:\allowbreak rank-\allowbreak two-\allowbreak height-\allowbreak schur-\allowbreak collapse & Higher\allowbreak Rank\allowbreak No\allowbreak Go.\allowbreak rank\allowbreak Two\allowbreak Height\allowbreak Schur\allowbreak Collapse \\
thm:\allowbreak gram-\allowbreak matrix-\allowbreak shadow & Higher\allowbreak Rank\allowbreak No\allowbreak Go.\allowbreak gram\allowbreak Matrix\allowbreak Shadow \\
thm:\allowbreak finite-\allowbreak window-\allowbreak no-\allowbreak go & Global\allowbreak Audit\allowbreak No\allowbreak Go.\allowbreak finite\allowbreak Window\allowbreak No\allowbreak Go \\
thm:\allowbreak finite-\allowbreak prime-\allowbreak cover-\allowbreak no-\allowbreak go & Support\allowbreak Prime\allowbreak No\allowbreak Go.\allowbreak finite\allowbreak Prime\allowbreak Cover\allowbreak No\allowbreak Go \\
def:\allowbreak support-\allowbreak prime-\allowbreak truncation-\allowbreak residual & Support\allowbreak Prime\allowbreak No\allowbreak Go.\allowbreak support\allowbreak Prime\allowbreak Truncation\allowbreak Residual \\
def:\allowbreak vsrc-\allowbreak algebra & Vsrc.\allowbreak vsrc\allowbreak Algebra \\
thm:\allowbreak vsrc-\allowbreak stage-\allowbreak i-\allowbreak consistency & Vsrc.\allowbreak vsrc\allowbreak Stage\allowbreak IConsistency \\
thm:\allowbreak vsrc-\allowbreak stage-\allowbreak ii-\allowbreak gz-\allowbreak partial & Vsrc.\allowbreak vsrc\allowbreak Stage\allowbreak IIGz\allowbreak Partial \\
thm:\allowbreak vsrc-\allowbreak stage-\allowbreak iii-\allowbreak translation & Vsrc.\allowbreak vsrc\allowbreak Stage\allowbreak IIITranslation \\
thm:\allowbreak kappa-\allowbreak r-\allowbreak normalization-\allowbreak equivalence & Kappa\allowbreak Normalization.\allowbreak kappa\allowbreak RNormalization\allowbreak Equivalence \\
\end{longtable}
\end{small}


% Per-axis bibliography; role-level citations are maintained in the prose.
\clearpage
\begingroup
\raggedright
\small
\bibliographystyle{plainnat}
\bibliography{references}
\endgroup
\end{document}
